\documentclass[11pt,reqno]{amsart}
\usepackage[T1]{fontenc}
\usepackage{lmodern,amsmath,amssymb,mathtools,microtype,booktabs,array,longtable}
\usepackage[margin=1in]{geometry}
\usepackage{graphicx}
\usepackage{xr-hyper}
\usepackage[hidelinks]{hyperref}
\ifdefined\pdfinfoomitdate\pdfinfoomitdate=1\fi
\ifdefined\pdftrailerid\pdftrailerid{}\fi
\ifdefined\pdfsuppressptexinfo\pdfsuppressptexinfo=15\fi
\numberwithin{equation}{section}
\newtheorem{theorem}{Theorem}[section]
\newtheorem{proposition}[theorem]{Proposition}
\newtheorem{lemma}[theorem]{Lemma}
\newtheorem{corollary}[theorem]{Corollary}
\theoremstyle{definition}\newtheorem{definition}[theorem]{Definition}
\newtheorem{example}[theorem]{Example}
\theoremstyle{remark}\newtheorem{remark}[theorem]{Remark}
\newcommand{\E}{\mathbb E}\newcommand{\Prb}{\mathbb P}
\newcommand{\R}{\mathbb R}\newcommand{\one}{\mathbf1}
\newcommand{\Q}{\mathcal Q}\newcommand{\TV}{\mathrm{TV}}
\newcommand{\cp}{\mathrm{cp}}\newcommand{\ev}{\mathrm{ev}}\newcommand{\st}{\mathrm{st}}
\DeclareMathOperator{\tr}{tr}\DeclareMathOperator{\dist}{dist}
\DeclareMathOperator{\law}{Law}\DeclareMathOperator{\rank}{rank}
\DeclareMathOperator{\supp}{supp}
\emergencystretch=2em

\externaldocument[][nocite]{supplement-refs}[General_Theta_Foundations_I_R70_Supplement.pdf]
\hypersetup{pdftitle={General Theta Foundations I: Acquired Geometry and Causal Resource Transfer},pdfauthor={Qian Qi}}
\title[Acquired geometry and causal resource transfer]{General Theta Foundations I:\\Acquired Geometry and Causal Resource Transfer}
\author{Qian Qi}\date{October 2026}
\subjclass[2020]{Primary 60G40, 62C20; Secondary 68Q45, 81P15, 90C22, 94A29}
\keywords{causal experiment, boundary optimality, stochastic realization, finite memory, robust prediction, quantum instrument, acquired geometry}
\begin{document}
\begin{abstract}
We study exact finite classical memory for prepared causal experiments whose observable response decays to zero. Every finite stochastic realization has an output-visible transient core and requires one additional live zero-output state. This reduction gives a sharp leakage law: for a single-frequency protected response, its least stochastic memory is one plus the smallest stochastic spectral dimension of the phase relative to the surviving population. Absolute leakage can be arbitrarily fast without reducing this count. A three-dimensional quantum channel is strictly contractive in trace norm and converges exponentially to one stationary state, yet its exact probability response has no finite classical realization. In one fixed squared-prediction embedding, finitely many rational physical models exclude every prescribed memory bound while the protected dimension remains six and the response rank remains four. Each fixed-bound risk gap is attained, positive and algebraically computable. Adding any positive background to the audit effect restores finite exact realization, with an explicit same-budget construction; its necessary memory diverges as the background vanishes. The inherited spectral, state-interior realization, stopping-rigidity and actual-history transfer theorems remain current with their separate hypotheses. Classical positive-realization and dominant-pole theory supply the underlying matrix methods; the present conclusions count the initial state, bounded decoder and physical stopping resources explicitly.
\end{abstract}
\maketitle
\section{Introduction and current theorem structure}\label{sec:current-overview}\label{sec:introduction}
Exact observer memory is a property of the performed task and its implementation interface, not merely the dimension or convergence rate of the protected physical system. The full-history optimum, its finite-memory approximation, and exact finite-memory attainment are distinct questions. Our current results concern a fourth distinction: leakage that reduces observable mass versus internal damping that changes phase relative to that mass.

Lemma~\ref{lem:r70-core} reduces every vanishing finite stochastic response to an output-visible transient matrix and proves an exact one-label increment. Proposition~\ref{prop:r70-relative-spectrum} extracts the observed pole in an irreducible block; it does not assume a false global Perron normalization of a reducible realization. Theorem~\ref{thm:r70-leakage} combines this necessity with an explicit construction to determine the exact leakage-dependent memory. All stochastic transition matrices, bounded randomized decoders and the prescribed initial label are included in both bounds.

Theorem~\ref{thm:r70-physical} realizes the obstruction in a strictly trace-contractive qutrit channel. Its complete probability response, not only a sign sequence, may need infinite classical memory even though the physical state converges exponentially. The fixed squared-loss embedding retains its actual audit and cost. Theorem~\ref{thm:r70-finite-tests} gives finite rational model witnesses of every prescribed memory obstruction on the same protected dimension. Their positive risk gaps are exact algebraic minima, not conclusions from numerical fitting. Theorem~\ref{thm:r70-background} supplies a positive counterpart: adding any strictly positive background to the audit effect restores finite exact stochastic realization, with all prefix and mixture labels counted.

Positive realizations, their dominant-pole restrictions and examples of small linear order requiring large positive order precede this work \cite{BenvenutiFarina04,BenvenutiFarina99,OCinneide90,Maier91}. Completely positive realizations and removal of unobservable positive components also have a developed theory \cite{MonrasWinter14}. We neither originate those methods nor equate a positive linear representation with a legal bounded-decoder observer. The proved contributions here are the exact constrained label increment, the leakage versus relative-damping classification, and the finite physical squared-task consequences. The prior spectral-memory theorem and its phase-sensitive rates are inherited and explicitly used, not claimed anew.

The instrument-wide interior theorem concerns a hidden simulator of complete instrument transcripts; its hidden state is not a free observation for a controller. The cost-rigidity theorem concerns the actual boundary of a single-command stopped experiment. The geometric theorem instead requires an actual-history converse and a common paid construction. These complementary results remain current, with their complete proofs in this article or its paired supplement. The retained ordered-measurement, reset-communication, rank-discrimination and nonintervening-sensing developments keep their independent provenance; none is inferred merely from the qutrit example.

Classical residual automata, finite linear rank, positive stochastic realization, dual witness support, protected quantum dimension, decoder width and program storage count different objects \cite{RabinScott59,ThonJaeger15,PiunovskiyZhang24}. The minimal reading guide isolates the new proof chain without making the entire preserved research archive a prerequisite.

\section{Prepared kernels, acquired task laws and observers}\label{sec:data}
All hidden-state, action and report spaces are standard Borel. A fixed unknown index $\theta\in\Theta$ specifies an actual preparation and positive normalized causal kernels. Positivity means nonnegative measures, not an everywhere strictly positive density. Dependence on the entire physical past may be included in the hidden state. A legal policy sees past visible reports and its admitted randomization, but not $\theta$ or a future audit. Iteration of these kernels defines its finite-prefix laws; the consistent countable path law is their kernel extension.

Each call has a cost $c_i\ge c_*>0$. The decision to make the next call is predictable. In the geometric route, the declared productive-call flag is also predictable; in the primitive-instrument route, productive reports are specified outcomes and can be random. Commitment is permitted only after a productive report, and occurs exactly once almost surely in each model. At that time a specified future audit $Z\in H$ is scored, where $H\subset\R^k$ is nonempty, compact and convex. The prediction is committed before that audit is observed. A prospective joint report--audit kernel is defined at every productive call, including rejected calls. The stopping coin is independent of the audit conditional on the entire current visible report and the past. Write $L=\operatorname{diam}(H)^2$; the case $L=0$ is trivial whenever a legal rule exists.

The core cost $C_\tau$ includes initial preparation charged to the core, failed calls and all productive calls through commitment. A separately executed final audit must be added to the raw budget. This convention is fixed before taking a risk infimum. Physical time, stored program constants and temporary computation are not identified with call count.

\subsection*{Three resource interfaces}
A checkpoint procedure may use its whole precommit history, but must output one element of a fixed dictionary of at most $M$ values. This counts its output width only; it is not a claim of $\log_2 M$ total working bits.

An autonomous observer has an alphabet $S$, $|S|\le M$, and one fixed parameter-independent program. At an intercall boundary its only persistent observer memory is its label $s\in S$. Its next command distribution is a fixed function of $s$. A physical port name, elapsed time, subprogram index, random seed or former report is not an extra free input. During a call its current command, current report and fresh coins are available. Whatever persists to another call must be encoded in $S$. A fixed decoder $r:S\to H$ reads only the committed label.

In the \emph{event-stopped} interface a post-report kernel emits $(s',e)\in S\times\{0,1\}$. An event $e=1$ ends the experiment and retains $s'$ for the decoder; no subsequent command is executed. This public commitment event is not a stored bit used in continuing computation. The same label may be used live on one path and terminal on another. In the \emph{state-stopped} interface a fixed map $h_*:S\to\{0,1\}$ determines halt status. The initial label is live, and type-one labels are globally absorbing terminal labels with no executed command rows. Entry into a terminal label is allowed only after a productive report. Thus live and terminal types are disjoint in this interface, and only in this interface by definition.

Write $R_i^\Theta(N,M)$, $i\in\{\cp,\ev,\st\}$, for the infimum of worst-case expected squared loss over programs with robust almost-sure commitment and $\sup_\theta\E_\theta C_\tau\le N$. An empty legal class has value $+\infty$. A program chosen once at random is not a free persistent seed. The ideal measurable class may have fixed real constants; effective implementation is stated separately in Section~\ref{sec:composition}.

\subsection*{A physically acquired comparison law}
Execute a declared prepared comparison experiment and obtain its joint report--audit probability $K_0(dx,dz)$ on $X\times H$. Its report marginal is $\nu$. The law is genuinely preparable and its specification has its own resource account. It need not be a member of $\Theta$. In particular no finite set of calibration data is silently promoted to a uniform certificate. Set
\begin{equation}\label{eq:derived}
 f(x)=\int zK_0(dz\mid x),\quad
 w(x)=\int\|z-f(x)\|^2K_0(dz\mid x),\quad
 d_r(x)=w(x)+\|f(x)-r\|^2.
\end{equation}
Disintegration gives Borel versions with $f\in H$; null sets may be completed arbitrarily. The profile is
\begin{equation}\label{eq:profile}
 \Q_m(b)=\min_{\substack{C\subset H,\ 1\le|C|\le m\\
                  0\le h\le b,\ \int h\,d\nu=1}}
        \int [w+\dist(f,C)^2]h\,d\nu,
 \qquad b\ge1,\ m\ge1.
\end{equation}
We use $+\infty$ for an infeasible cap or width. The quantities $f,w$ are task coordinates of acquired predictive information, not a complete substitute for the predictive quotient. Their derivation and the quotient construction in Appendix~\ref{sec:quotient} precede any use of local dimension.

\subsection*{Lower certificates on actual histories}\label{sec:lower-cert}
Fix \emph{one} actual model $\theta_*\in\Theta$, independent of the competing program. At each actual reached productive pre-call history $h$ and legal command $a$, write $P_{h,a}(dy,dz)$ for its conditional full current report--audit law. Let $P^0_{h,a}$ be a measurable comparison row on the same report space, with a Borel mark $x=\pi_{h,a}(y)$ of law $\nu$ and
\begin{equation}\label{eq:full-moments}
 \E^0[Z\mid y]=f(\pi_{h,a}y),\qquad
 \E^0[\|Z\|^2\mid y]=w(\pi_{h,a}y)+\|f(\pi_{h,a}y)\|^2.
\end{equation}
These are conditional identities given the \emph{full} report, not just its mark. Versions and inequalities are required on actual reachable histories, not only almost everywhere under an unrelated comparison path law.

Fix $0<l_-\le1\le u_-<\infty$. For all Borel $q:Y\to[0,1]$ and all finite-image Borel $r:Y\to H$, suppose predictable nonnegative defects $a_-(h,a),e_-(h,a)$ obey
\begin{align}
 \int q\,d(P_{h,a})_Y
 &\le u_-\int q\,d(P^0_{h,a})_Y+a_-(h,a),\label{eq:lower-mass}\\
 \int q(y)\|z-r(y)\|^2P_{h,a}(dy,dz)
 &\ge l_-\int q(y)\|z-r(y)\|^2P^0_{h,a}(dy,dz)-e_-(h,a).
 \label{eq:lower-task}
\end{align}
Independent current randomization can be adjoined to both rows: integrating inequalities for each fixed random value proves the same inequalities with the same defects. All command-conditioned lower inequalities hold before averaging the chosen command.

Let $A_\tau$ count productive calls. At this same actual model require, for every legal competing procedure,
\begin{equation}\label{eq:effort}
 \E_* C_\tau\ge\kappa_-\E_* A_\tau,\qquad\kappa_->0.
\end{equation}
At budget $N$ suppose, uniformly over all robustly feasible competitors,
\begin{equation}\label{eq:occupation}
 \E_*\sum_{j\le\tau:\ j\ {\rm productive}}a_-(H_{j-1},a_j)\le A_N,
 \qquad
 \E_*\sum_{j\le\tau:\ j\ {\rm productive}}e_-(H_{j-1},a_j)\le E_N.
\end{equation}
These are actual occupation certificates, not properties merely of the proposed upper program. Uniform row defects give $A_N\le(N/\kappa_-)\bar a_-$ and $E_N\le(N/\kappa_-)\bar e_-$. Lemma~\ref{lem:drift} supplies a nonuniform alternative.

\subsection*{One common upper construction}\label{sec:upper-cert}
Independently, let $\mathcal G$ be a parameter-independent program which, from each of its start histories, reaches one productive report. Its conditional expected attempt cost is at most $\kappa_+>0$ on every surviving history in every $\theta\in\Theta$. It uses $B$ live roles and returns to its start label on rejection; no reset of the physical state is assumed. It forms a virtual mark $x$, a current choice $J$ and a permitted export $T_Jr$ of its task readout. For every finite dictionary $C$ assume
\begin{equation}\label{eq:exports}
 |\{T_jr:j\in J,\ r\in C\}|\le D_0|C|,
 \qquad D_0\in\mathbb N.
\end{equation}
Every continuing phase and coin needed for this construction is included in the $B$ live roles. For fixed $0<l_+\le1\le u_+<\infty$ and $a_+,e_+\ge0$, assume for all Borel $q:X\to[0,1]$ and finite-image Borel $r:X\to H$, conditionally on every such history,
\begin{align}
 \E[q(x)]&\ge l_+\int q\,d\nu-a_+,\label{eq:upper-mass}\\
 \E[q(x)\|Z-T_Jr(x)\|^2]
 &\le u_+\int q(x)d_{r(x)}(x)\nu(dx)+e_+.
 \label{eq:upper-task}
\end{align}
The expectation here includes the within-attempt choice $J$. Unlike the lower, the upper is not required to hold for each choice separately. The choices retained after this call are nevertheless charged through the actual export or live role. These hypotheses express one common causal construction, not modelwise optimizers.

\section{Transient output faces and leakage-sensitive memory}\label{sec:r70-leakage}
Uniform decay of a physical state and decay of its observable probabilities do not measure the same obstruction to exact stochastic realization. In this section leakage removes population from a protected subspace, whereas internal depolarization changes the spectrum relative to that surviving population. We prove that the first operation costs exactly one additional live label in the nondegenerate single-frequency family, independently of how rapid the leakage is. We then obtain finite-model witnesses on a fixed six-dimensional squared-prediction experiment, and an exact finite-realization result when a positive background is added to the terminal effect.

Nonnegative realizations, dominant-pole obstructions and transient-state reductions have classical antecedents \cite{BenvenutiFarina04,OCinneide90,Maier91,MonrasWinter14}. We do not claim to originate these methods. The statements below specify the constrained stochastic decoder, prescribed initial label, sharp label increment, physical leakage channel and finite squared-task tests. All necessary positive-matrix arguments are included.

\subsection{The output-visible transient core}
For a nonzero sequence $u=(u_n)_{n\ge0}$ in $[0,1]$ tending to zero, let $s(u)$ be its least stochastic realization size:
\[
 u_n=e_1^{\mathsf T}Q^nd,\qquad Q\ge0,\quad Q\one=\one,\quad 0\le d\le\one.
\]
The initial label is fixed, not a freely randomized initial distribution. Define $t(u)$ by the same formula with a nonnegative row-substochastic matrix $S$, $S\one\le\one$, whose spectral radius is strictly below one. Both minima are infinite if no finite realization exists.

\begin{lemma}[Exact transient-core reduction]\label{lem:r70-core}
For every such sequence,
\begin{equation}\label{eq:r70-core-count}
 s(u)=t(u)+1.
\end{equation}
More precisely, any $M$-label stochastic realization has a nonempty set of at most $M-1$ reachable labels that can eventually give positive output. Its restriction $S$ realizes $u$, is transient and uses the same initial label. If also $u_n\le C\lambda^n$ for $0<\lambda<1$, then the spectral radius of this restriction is at most $\lambda$.
\end{lemma}
\begin{proof}
Remove labels not reachable from label one. Call a remaining label active if a positive-probability path from it reaches some label $j$ with $d_j>0$. The initial label is active because $u$ is nonzero: there is some $n$ with $e_1^{\mathsf T}Q^nd>0$, so a positive-weight path from label one reaches a label with positive decoder output. This remains true when $u_0=0$. A path that leaves the active set cannot return: a return would make every intervening reachable label active. Thus, with $S$ the active principal submatrix and $d_A$ its decoder, $u_n=e_1^{\mathsf T}S^nd_A$.

For each active $i$, choose a path of length $a_i$ from the start to $i$ of weight $p_i>0$. For each active $j$, choose a path of length $b_j$ from $j$ with positive terminal expected output $r_j>0$. Positivity gives, for every $n$,
\begin{equation}\label{eq:r70-path-bound}
 p_ir_j(S^n)_{ij}\le u_{a_i+n+b_j}.
\end{equation}
Consequently every entry of $S^n$ tends to zero. In finite dimension this forces every eigenvalue to have modulus less than one. In the exponential case it gives $(S^n)_{ij}\le C_{ij}\lambda^n$, and applying this bound to an eigenvector proves the stronger spectral-radius bound. There must be a reachable inactive label: otherwise $S$ would be stochastic and have eigenvalue one. Hence the active size is at most $M-1$, proving $t(u)+1\le s(u)$.

Conversely a $k$-label transient realization extends to the stochastic matrix and decoder
\begin{equation}\label{eq:r70-cemetery}
 \widehat Q=\begin{pmatrix}S&\one-S\one\\0&1\end{pmatrix},
 \qquad \widehat d=\binom{d_A}{0}.
\end{equation}
They realize the same sequence from the same initial label with $k+1$ labels. This proves the reverse inequality, including the infinite case. The added label is a live zero-output label until the actual herald; it is not permission to stop at an unobserved physical leakage event.
\end{proof}

We recall a small positive-matrix fact with its proof. If $B\ge0$ is irreducible and has a nonzero eigenvalue, it has a vector $h>0$ with $Bh=Rh$, where $R$ is its spectral radius. Indeed, for $B+\epsilon\one\one^{\mathsf T}$ the map sending a simplex vector to its normalized image has a fixed point by the finite-dimensional simplex fixed-point theorem. The fixed point is strictly positive. Along $\epsilon\downarrow0$ a convergent subsequence gives $Bh=Rh$, $h\ge0$, $\sum_i h_i=1$. If a coordinate vanished, the eigenvector equation and irreducibility would propagate that zero to every coordinate. Thus $h>0$, and $R>0$. The similarity $R^{-1}\operatorname{diag}(h)^{-1}B\operatorname{diag}(h)$ is stochastic, so its maximum-row-sum bound shows that $R$ is indeed the spectral radius. This is the Perron normalization used below, not a new Perron theorem.

\begin{proposition}[A relative spectral obstruction for decaying outputs]\label{prop:r70-relative-spectrum}
Suppose an $M$-label stochastic response satisfies $0\le u_n\le C\lambda^n$ with $0<\lambda<1$. If its rational generating function has a nonreal pole at $\zeta^{-1}$, then
\begin{equation}\label{eq:r70-relative-spectrum}
 \zeta/\lambda\in\mathcal K_{M-1}.
\end{equation}
Here $\mathcal K_j$ is the stochastic eigenvalue region defined in Section~\ref{sec:r69-spectral}. It is the decay-normalized mode, not merely the absolute mode $\zeta$, that constrains the active state count.
\end{proposition}
\begin{proof}
The core in Lemma~\ref{lem:r70-core} gives
$\sum_nu_nz^n=e_1^{\mathsf T}(I-zS)^{-1}d_A$. Hence $\zeta$ is an eigenvalue of $S$. Order its strongly connected components to obtain a block triangular matrix. One irreducible diagonal block $B$ has that eigenvalue. Its spectral radius $R$ obeys $|\zeta|\le R\le\lambda$. Perron normalization gives a stochastic matrix $T$ of size $l\le M-1$ with eigenvalue $\zeta/R$ and eigenvector $v$.

A stationary probability row $\pi$ for $T$ exists: take a convergent subsequence of the averages of any probability row and its first $n-1$ iterates. Since $\zeta/R\ne1$, one has $\pi v=0$. Put $a=R/\lambda\in(0,1]$. The stochastic matrix
$T'=aT+(1-a)\one\pi$ satisfies $T'v=(\zeta/\lambda)v$. Padding by unreachable self-loop labels, if needed, proves membership in $\mathcal K_{M-1}$. No global Perron scaling of the reducible core has been assumed.
\end{proof}

The last qualification matters. A reducible matrix such as
$S=\left(\begin{smallmatrix}1/2&1/4\\0&1/4\end{smallmatrix}\right)$
has spectral radius $1/2$, but no strictly positive $h$ satisfies $Sh\le h/2$. The proof normalizes the irreducible block carrying the observed pole, rather than silently imposing that false global inequality.

\subsection{An exact leakage law for every damping and phase}
Fix $0<\lambda<1$, $0<\alpha\le1$, $0<\gamma\le1$, and $\sin\vartheta\ne0$. Let
\begin{equation}\label{eq:r70-leaky-response}
 u_n=\frac{\lambda^n}{2}
       [1+\gamma\alpha^n\cos(n\vartheta+\varphi)].
\end{equation}
The scalar $\lambda$ is survival against leakage; $\alpha$ is coherence retention within the surviving qubit. These parameters are known properties of the prepared target channel, not extra observations available to the observer.

\begin{theorem}[Exact leakage increment]\label{thm:r70-leakage}
For the nondegenerate family \eqref{eq:r70-leaky-response},
\begin{equation}\label{eq:r70-exact}
 s(u)=1+m(\alpha,\vartheta)
     =1+\min\{k:\alpha e^{i\vartheta}\in\mathcal K_k\}.
\end{equation}
This holds over all stochastic matrices and randomized decoder rows, with a prescribed initial label. In particular the value is independent of every $0<\lambda<1$, every nonzero contrast and every preparation phase. For $\alpha<1$ it is finite. At $\alpha=1$ it is $h+1$ if $e^{i\vartheta}$ has finite order $h$, and it is infinite otherwise.

The full two-output Hankel response has rank four, regardless of this possibly much larger or infinite exact memory. In the fixed squared-task embedding it is the least live memory at the original cost $1/p$, with event-stopped and state-stopped total-label upper counts $\max\{s(u),2\}$ and $s(u)+2$ when finite.
\end{theorem}
\begin{proof}
The three modes $\lambda$, $\lambda\alpha e^{i\vartheta}$ and $\lambda\alpha e^{-i\vartheta}$ in $u$ are distinct and have nonzero coefficients because $\lambda>0$, $\alpha>0$, $\gamma>0$ and $\sin\vartheta\ne0$. The phase factor $e^{i\varphi}$ never vanishes, including at initial-output zeros. Thus its generating function has a pole at $(\lambda\alpha e^{i\vartheta})^{-1}$. Proposition~\ref{prop:r70-relative-spectrum} makes every finite realizing size $M$ satisfy $\alpha e^{i\vartheta}\in\mathcal K_{M-1}$. Theorem~\ref{thm:r69-spectral} proves the required lower and the infinite case.

For the upper take a minimal realization $(Q,d)$ of the unattenuated response
$b_n=[1+\gamma\alpha^n\cos(n\vartheta+\varphi)]/2$.
Replace its continuing block by $\lambda Q$ and append a live zero-output label, as in \eqref{eq:r70-cemetery}. Its output is exactly $\lambda^nb_n$ from time zero. This proves equality without an extra initialization operation. The unit-circle alternatives follow from the already proved root-of-unity characterization.

The scalar response spans its three exponential modes. Its complement adds the independent constant mode one, because $\lambda<1$. Vandermonde independence therefore gives full response rank four. Theorem~\ref{thm:r67-embedding} forces this entire probability response in the actual squared task; the role-packing statement is then inherited without identifying protected physical states with live observer labels.
\end{proof}

\begin{corollary}[Noise laws depend on relative damping]\label{cor:r70-rates}
For fixed $0<\lambda<1$ and $0<\alpha<1$, the worst-phase exact count is
\[
 \max_{\sin\vartheta\ne0}s(u)
 =1+\max\{3,\lceil\pi/\arccos\alpha\rceil\}.
\]
For every fixed phase with the global badly-approximable certificate of Theorem~\ref{thm:r69-cubic}, $s(u)-1$ has the same matching order $(1-\alpha)^{-1/3}$. Neither rate depends on leakage $\lambda$, even when leakage is allowed to vary simultaneously with $\alpha$. At the original algebraic phase, the first nontrivial threshold is exactly four live labels for $0<\alpha\le5/(3+4\sqrt3)$, and at least five beyond it. No badly-approximable claim is made for that algebraic phase.
\end{corollary}
\begin{proof}
Substitute \eqref{eq:r70-exact} into the exact worst-phase, fixed-phase and three-label results. All subtractions concern live labels, including the necessary zero-output label; no terminal label has been subtracted from an event-stopped dictionary.
\end{proof}

\subsection{A strictly contractive quantum channel with infinite exact memory}
Let $K=\mathbb C^2\oplus\mathbb C|\bot\rangle$, let $P$ project onto the qubit and let $\Omega=|\bot\rangle\langle\bot|$. With the qubit channel $\Phi_\alpha$ and preparation from Section~\ref{sec:r69-spectral}, define
\begin{equation}\label{eq:r70-channel}
 \Psi_{\lambda,\alpha}(X)
 =\lambda\Phi_\alpha(PXP)
  +[\tr((I-P)X)+(1-\lambda)\tr(PX)]\Omega.
\end{equation}
The qubit output is embedded in its original block and cross-block coherences are discarded. Prepare the qubit and audit the qubit effect $E=(I+\sigma_z)/2$, extended by zero at $\bot$.

\begin{theorem}[Uniform physical mixing versus exact response memory]\label{thm:r70-physical}
Equation \eqref{eq:r70-channel} is a completely positive trace-preserving channel with unique stationary state $\Omega$. For all density matrices $\rho,\sigma$ and all $n\ge0$,
\[
 \|\Psi_{\lambda,\alpha}(\rho)-\Psi_{\lambda,\alpha}(\sigma)\|_1
 \le\lambda\|\rho-\sigma\|_1,
 \qquad
 \|\Psi_{\lambda,\alpha}^n(\rho)-\Omega\|_1\le2\lambda^n.
\]
Its prepared audit response is exactly \eqref{eq:r70-leaky-response}, so its exact stochastic observer memory is \eqref{eq:r70-exact}. In particular, at $\alpha=1$ and the algebraic rotation $(3+4i)/5$, every $0<\lambda<1$ gives infinite exact memory despite strict trace-norm contraction and exponential convergence to one stationary state. For algebraic $\lambda$ and algebraic preparation data, all primitive superoperator entries are algebraic.

The squared-prediction embedding has protected dimension six, fixed dictionary $\{-1,1\}$, three actual audit outcomes and least expected cost $1/p$. In the specified irrational, internally undamped case, every finite live-memory class has an attained strictly larger risk than the full-history value $2$. A prefix with $T$ exact-output labels and one live zero-output tail label has the same cost and excess risk at most
\begin{equation}\label{eq:r70-approx}
 \frac{2\varepsilon p(q\bar t\lambda)^T}{1-q\bar t\lambda},
 \qquad q=1-p,\quad\varepsilon=1/4.
\end{equation}
\end{theorem}
\begin{proof}
Write $\Xi(X)=\Phi_\alpha(PXP)+\tr((I-P)X)\Omega$ and
$\mathcal E(X)=\tr(X)\Omega$. Both are channels: compression, the declared qubit channel and scalar replacement are completely positive, and their output traces sum to the input trace. Then $\Psi_{\lambda,\alpha}=\lambda\Xi+(1-\lambda)\mathcal E$. On a traceless Hermitian difference the second term vanishes. Trace-norm contraction of a channel follows directly by applying positivity and trace preservation to the positive and negative parts of that difference. This proves the contraction inequality. Since $\Omega$ is fixed, iteration gives the displayed bound and uniqueness.

Starting in the qubit, its surviving block after $n$ calls is $\lambda^n\Phi_\alpha^n(\rho)$; the complementary mass is at $\bot$. Taking the effect gives \eqref{eq:r70-leaky-response}. The preceding theorem proves the memory claim. At $(3+4i)/5$ the unit phase is not a root of unity, by the rational algebraic-integer argument already proved. The spectral obstruction is in the decaying output-visible block, not at the peripheral spectrum of the entire trace-preserving channel.

The embedding of Theorem~\ref{thm:r67-embedding} doubles the protected dimension and retains its real squared audit. Its unique optimal probability kernel must reproduce $u$. Compactness of each fixed-memory controller class and the uniform geometric word tail give attainment; equality at a finite size would contradict the exact-memory theorem. For the approximation, store $u_0,\ldots,u_{T-1}$ and thereafter output probability zero. Since $u_n\le\lambda^n$, summing the exact risk identity \eqref{eq:r67-generating-risk} yields \eqref{eq:r70-approx}. Nothing in this construction observes leakage or changes first-herald stopping.
\end{proof}

This is not the earlier unique-sign counterexample under depolarization. Here the required response is probability valued and its absolute nonconstant eigenvalues all lie strictly inside the unit disk. The impossibility comes from their phases relative to the visible decay rate. The stationary state is on a face, so the interior hypothesis of the instrument-wide realization theorem is genuinely unavailable. With $\gamma=1$ and $\varphi=\pi$, moreover, $u_0=0$ and $0\le u_n\le\lambda$ for all $n\ge1$. The one-label zero response is uniformly within $\lambda$ of this law while exact memory remains infinite for every positive $\lambda$. Exact realization is not a robust approximate-memory assertion.

\subsection{Finitely many physical models detect the obstruction}
The continuum of herald-flag parameters used to force an arbitrary full response can be replaced, against a prescribed finite-memory bound, by a finite set whose cardinality depends only on that bound and the target's linear order.

\begin{theorem}[Finite rational-resolvent tests and fixed-dimensional witnesses]\label{thm:r70-finite-tests}
Suppose a target channel response has generating function $U(z)=P(z)/D(z)$ with
$\deg P\le r-1$, $\deg D\le r$, and $D(z)\ne0$ on $[0,q\bar t]$. Fix $M\ge1$ and any $M+r$ distinct parameters $t_j\in[0,\bar t]$. In the squared embedding, an observer with at most $M$ live labels has risk exactly $2$ on the finite family $(t_j,\sigma)$, $\sigma\in\{-1,1\}$, if and only if it realizes the entire target response at every failure count.

Consequently for the algebraic leaking qubit with $\lambda=1/2$, $\alpha=\gamma=1$, $e^{i\vartheta}=(3+4i)/5$ and $\varphi=0$, the following holds for every $M$: a six-dimensional prepared squared experiment with only $2(M+3)$ rationally specified models, one unit-cost command and the fixed two-value dictionary has boundary value $2$, but every at-most-$M$-label observer has risk at least $2+\delta_M$, where $\delta_M>0$ is an attained, computable real algebraic number. One may take $p=1/2$, $\bar t=1/2$ and $t_j=j/[2(M+2)]$, $0\le j\le M+2$. Neither the physical dimension nor the target response order increases with $M$.
\end{theorem}
\begin{proof}
Boundary feasibility still forces first-herald commitment in this finite model family, since its physical herald probability is always $p$. An $M$-label rule has
$B(z)=e_1^{\mathsf T}(I-zQ)^{-1}d=P_Q(z)/D_Q(z)$,
where $\deg P_Q\le M-1$, $\deg D_Q\le M$ and $D_Q(z)\ne0$ for $0\le z<1$. Smaller matrices may be padded by unreachable rows.

The two signs and \eqref{eq:r67-generating-risk} make risk $2$ equivalent to $B(qt_j)=U(qt_j)$ at every chosen parameter. The polynomial
$P_QD-PD_Q$ has degree at most $M+r-1$ and vanishes at $M+r$ distinct points. It is therefore zero. Equality of the power series about zero gives every response coefficient. The converse follows immediately from the same risk identity. This is rational interpolation with known degree, not extrapolation from arbitrary fitted samples.

For the stated target,
\[
 U(z)=\frac1{2(1-\lambda z)}+
       \frac{1-\lambda cz}{2(1-2\lambda cz+\lambda^2z^2)},
 \qquad c=3/5,
\]
so $r=3$ suffices and all coefficients are rational. The exact leakage theorem excludes every finite realization. Hence the finite tests exclude equality for the specified $M$. The full-history target kernel still attains $2$ in every selected model. Compact controller simplexes and the geometric tail imply the fixed-$M$ minimum is attained, so its gap $\delta_M$ is strictly positive.

Finally introduce real variables $Q,d$ and vectors $v_j$ satisfying
$(I-qt_jQ)v_j=d$. These matrices are invertible for every stochastic $Q$, and the achieved risk is
$2+2\varepsilon p\max_j|(v_j)_1-U(qt_j)|$.
All defining constraints are polynomial over the rationals. Real-algebraic elimination computes the attained minimum and an algebraic minimizing controller; in particular the positive gap is algebraic and computable. This is an exact finite procedure, not a polynomial-time claim or a numerical certificate inferred from local optimization. The general procedure is proved here; a full elimination at every $M$ has not been executed in the finite diagnostics.
\end{proof}

Thus finite linear rank four does not bound the classical memory forced by a finite family of actual squared-loss tests, even on a fixed protected dimension. The number of models in this conclusion grows with the tested memory bound. It is not a claim that one fixed finite model list excludes every finite observer.

\subsection{Positive background regularizes the output face}
The leakage obstruction is sensitive to the face on which the terminal effect vanishes. Replace the effect by
$E_\eta=\eta I/2+(1-\eta)E$, $0<\eta<1$, on the same qutrit. The response becomes
\[
 u_n^{(\eta)}=\eta/2+(1-\eta)u_n.
\]
This changes a real terminal measurement, not the protected channel or the observer's information.

\begin{theorem}[Exact finite realization after background regularization]\label{thm:r70-background}
For all parameters of Theorem~\ref{thm:r70-leakage} and every $0<\eta<1$, the response $u^{(\eta)}$ has a finite exact stochastic realization. If $T\ge1$ satisfies
\begin{equation}\label{eq:r70-background-depth}
 (1-\eta)\lambda^T\le\eta/2,
\end{equation}
then at most $T+2+m(\lambda\alpha,\vartheta)$ live labels suffice. The construction uses a prescribed initial label and attains the squared-task optimum at the same cost $1/p$.

For fixed $\lambda\in(0,1)$, $\alpha=1$ and irrational phase, the least exact size tends to infinity as $\eta\downarrow0$, while the displayed upper is $O_\lambda(\log(1/\eta))$. No matching lower rate in $\eta$ is asserted. For algebraic nondegenerate inputs this finite upper and the inherited fixed-size response test give a terminating search for a minimal exact controller.
\end{theorem}
\begin{proof}
Put $a=(1-\eta)\lambda^T/2$ and $\psi=T\vartheta+\varphi$. The suffix starting at $T$ is the average of the two probability responses
\[
 f_k=\eta/2+2a\lambda^k,\qquad
 g_k=\eta/2+2a\gamma\alpha^T(\lambda\alpha)^k\cos(k\vartheta+\psi).
\]
The first has a two-label transient/absorbing realization with decoder values $\eta/2+2a$ and $\eta/2$. For the second, use the normalized eigenvector construction of Theorem~\ref{thm:r69-spectral} at eigenvalue $\lambda\alpha e^{i\vartheta}$, with decoder $\eta/2+2a\gamma\alpha^T\operatorname{Re}(e^{i\psi}v_j)$. Condition \eqref{eq:r70-background-depth} gives $2a\le\eta/2$, so all these decoder values lie in $[0,\eta]\subset[0,1]$.

Store a deterministic chain for $u_0^{(\eta)},\ldots,u_{T-1}^{(\eta)}$. At its final failure transition enter the two disjoint suffix machines with weights $1/2,1/2$. No additional initialization label is needed: the last prefix row makes and stores the choice within the counted blocks. At every later failure count the output is $(f_k+g_k)/2=u_{T+k}^{(\eta)}$. This proves the exact count and preserves the physical stopping rule.

If exact sizes stayed bounded along $\eta\downarrow0$, pad their stochastic matrices and decoders to one size and take a convergent subsequence in the compact row simplexes. Every finite coefficient is continuous in those variables. The limit would realize $u$ at $\eta=0$, contradicting Theorem~\ref{thm:r70-leakage}. The logarithmic upper follows by solving \eqref{eq:r70-background-depth} with fixed $\lambda$. For algebraic input, the target has a supplied linear representation of dimension at most four. The finite word-equality criterion of Proposition~\ref{prop:finite-realization-test}, together with this size upper, terminates a minimal exact-size search. Algebraic controller encodings and exact sampling resources remain distinct from label count.
\end{proof}

\section{Spectral and arithmetic laws of exact response memory}\label{sec:r69-spectral}
The instrument realization theorem preserves every terminal effect, whereas a prescribed task can require substantially less memory. We now determine the exact memory problem for a whole family of single-frequency protected responses, not just one damping factor or the previously solved three-label case. The criterion reduces the problem to classical stochastic eigenvalue regions. A self-contained sector argument and elementary rational approximation then give an exact worst-phase law and matching growth at every fixed badly approximable phase. None of the spectral-region, positive-realization, or sector machinery is claimed as newly invented \cite{DmitrievDynkin46,JohnsonPaparella17,KirklandSmigoc22}. The contributions considered here are their exact response-level equivalence with a prescribed initial label, their consequences for the physical squared task, and the phase-sensitive matched resource law.

\subsection{A physical family and the exact spectral criterion}
Let $0<\alpha\le1$, $0<\gamma\le1$, and $\vartheta,\varphi\in\mathbb R$, with $\sin\vartheta\ne0$. Consider the binary response
\begin{equation}\label{eq:r69-response}
 b_n=\frac12\bigl[1+\gamma\alpha^n\cos(n\vartheta+\varphi)\bigr],
 \qquad n\ge0.
\end{equation}
The rotation, contrast and preparation phase specify one prepared instance; independence of the minimum size does not assert one fixed decoder for unknown phases or contrasts. It is physically prepared on a qubit: take the Bloch vector
$\gamma(\sin\varphi,0,\cos\varphi)$, the channel
\[
 \Phi(X)=\alpha U_\vartheta XU_\vartheta^*
               +(1-\alpha)\tr(X)I/2,\qquad
 U_\vartheta=\begin{pmatrix}\cos(\vartheta/2)&-\sin(\vartheta/2)\\
                  \sin(\vartheta/2)&\cos(\vartheta/2)\end{pmatrix},
\]
and terminal effect $E=(I+\sigma_z)/2$. Its Born probabilities are exactly \eqref{eq:r69-response}. Theorem~\ref{thm:r67-embedding} embeds this response into the same three-outcome squared-loss experiment, with protected dimension four, fixed decoder dictionary $\{-1,1\}$, and least expected cost $1/p$. It forces the probability of output $+1$ after $n$ failures to equal $b_n$. Thus the mathematical realization problem below is the exact optimal-observer problem of an actual bounded squared task, not a sequence supplied by a loss oracle.

Write $m(\alpha,\vartheta;\gamma,\varphi)$ for the least $M$ for which
\[
 b_n=e_1^{\mathsf T}Q^nd\quad(n\ge0),\qquad
 Q\ge0,\quad Q\mathbf1=\mathbf1,\quad d\in[0,1]^M.
\]
The initial label is prescribed and counted. Let
\[
 \mathcal K_M=\{z\in\mathbb C:z\text{ is an eigenvalue of some }
                         M\times M\text{ row-stochastic matrix}\}.
\]
These are the classical Karpelevi\v c regions; their full boundary classification is prior work \cite{JohnsonPaparella17}. We will not use that classification as an unproved step in any bound below.

\begin{theorem}[Exact spectral realization, with no extra initial label]\label{thm:r69-spectral}
For the nondegenerate family \eqref{eq:r69-response},
\begin{equation}\label{eq:r69-minimal}
 m(\alpha,\vartheta;\gamma,\varphi)
       =\min\{M:\alpha e^{i\vartheta}\in\mathcal K_M\}.
\end{equation}
The value is independent of every nonzero contrast $\gamma$ and of the preparation phase $\varphi$; denote it by $m(\alpha,\vartheta)$. The criterion is necessary and sufficient over all stochastic matrices, including reducible matrices and arbitrary randomized decoders. For $0<\alpha<1$ it is finite. For $\alpha=1$ and $\vartheta/(2\pi)=a/q$ in lowest terms it is exactly $q$; for irrational $\vartheta/(2\pi)$ it is infinite.

In the squared-task embedding, these are exact live-label counts at cost $1/p$. At each fixed smaller memory bound, the attained risk is strictly above the full-history value $2$. Event-stopped and state-stopped total-label upper counts remain respectively $\max\{M,2\}$ and $M+2$.
\end{theorem}
\begin{proof}
Put $z=\alpha e^{i\vartheta}$. The response has the three distinct exponential modes $1,z,\bar z$, with nonzero coefficients. Cayley--Hamilton for any realizing $Q$ makes its characteristic polynomial an annihilating shift polynomial of $b_n$. Substitution and Vandermonde independence force that polynomial to vanish at $z$. This proves necessity without assuming a minimal linear representation has positive coordinates.

Conversely let $Qv=zv$ with $v\ne0$. Normalize a largest-modulus coordinate to one by multiplying $v$ by one complex scalar, and permute that coordinate into position one. Then $v_1=1$ and $|v_j|\le1$. The column
\[
 d_j=\tfrac12[1+\gamma\operatorname{Re}(e^{i\varphi}v_j)]
\]
lies in $[0,1]^M$. Row normalization and $Q^nv=z^nv$ give $e_1^{\mathsf T}Q^nd=b_n$ for every $n$. The same $M$ labels suffice: there is no preliminary random initialization or extra anchor. Unreachable labels can be pruned. Finiteness for $\alpha<1$ follows from the regular-polygon construction proved below, or from the inherited polygon upper in Theorem~\ref{thm:qubit-positive-phase}.

At $|z|=1$, equality in the triangle inequality at a largest-modulus eigenvector coordinate forces each positive successor to have the same modulus and value $z$ times the current value. Following positive successors reaches a directed cycle of length $l\le M$, so $z^l=1$. An irrational phase is therefore impossible. If $z$ has order $q$, every such cycle has length divisible by $q$, giving $M\ge q$; a $q$-cycle realizes it. Finally the squared embedding forces the whole response. Compact attainment at fixed $M$ was proved in Proposition~\ref{prop:r67-density}; equality of its risk with $2$ would give the excluded exact response. The existing packing lemma gives the terminal counts.
\end{proof}

The nondegeneracy assumptions cannot be suppressed. Zero contrast gives one constant label. At $\alpha=0$, $b_0=(1+\gamma\cos\varphi)/2$ and $b_n=1/2$ for $n\ge1$; one label suffices if $\gamma\cos\varphi=0$, and otherwise two are necessary and sufficient. Real phases and constant cases are classified in the supplement. In particular a vanishing oscillation amplitude is not the limit of an exact-memory lower valid for nonzero amplitude.

For algebraic $x=\operatorname{Re}z$, $y=\operatorname{Im}z\ne0$, spectral membership has the compact polynomial certificate
\begin{equation}\label{eq:r69-bilinear}
 \begin{gathered}
 Q\ge0,\ Q\mathbf1=\mathbf1,\quad
 Qa=xa-yc,\quad Qc=ya+xc,\\
 a_1=1,\ c_1=0,\quad a_j^2+c_j^2\le1\quad(1\le j\le M).
 \end{gathered}
\end{equation}
It has $M^2+2M$ real variables and polynomial degree at most two. Normalization of an eigenvector proves equivalence in both directions. Real-algebraic elimination decides it at each fixed $M$, and the finite polygon upper terminates a minimal-size search for $|z|<1$. This improves the form of the certificate, not the general complexity of real-algebraic elimination. A chosen algebraic preparation and contrast then yield algebraic decoder entries. State labels, encodings of those entries, exact-sampling resources and runtime remain separate accounts.

\subsection{The sector bound, its equality case, and exact finite phases}
We include the classical stochastic sector argument to expose every load-bearing inequality. The equality structure is consistent with the Type-0 classification of Kirkland and \v Smigoc \cite{KirklandSmigoc22}; it is not a new classification of stochastic spectral arcs.

\begin{lemma}[Stochastic sector inequality and rigidity]\label{lem:r69-sector}
If $z=x+iy$, $y\ne0$, is an eigenvalue of an $M\times M$ stochastic matrix, then $M\ge3$ and
\begin{equation}\label{eq:r69-sector}
 |y|\le(1-x)\cot(\pi/M).
\end{equation}
If equality holds, then, up to simultaneous relabelling and reversing the orientation when needed, that matrix is
\begin{equation}\label{eq:r69-lazy-cycle}
 Q=(1-t)I+tC_M,\qquad 0<t\le1,
\end{equation}
where $C_M$ is the directed $M$-cycle. Its eigenvalue on the equality ray is $1-t+te^{2\pi i/M}$, or its conjugate.
\end{lemma}
\begin{proof}
Conjugate the eigenvector when necessary so $y>0$. Let $P$ be the convex hull of its complex coordinates. The matrix
$e^{s(Q-I)}=e^{-s}\sum_{n\ge0}s^nQ^n/n!$ is stochastic for $s\ge0$, so multiplication by $e^{s(z-1)}$ maps $P$ into itself. Since $z$ is nonreal, $1-x>0$, and the spiral of any nonzero coordinate visits every angle, a finite set of points of this spiral surrounds zero. Hence zero lies in the interior of $P$.

List its $l\le M$ vertices counterclockwise as $r_je^{i\theta_j}$, with $r_j>0$. Their consecutive angular gaps $\Delta_j$ belong to $(0,\pi)$ and sum to $2\pi$. Put $u=1-x>0$ and $v=y>0$. The initial velocity $(-u+iv)r_je^{i\theta_j}$ must point into the half-plane to the left of the outgoing edge. Taking its signed two-dimensional cross product with that edge gives
\[
 r_j/r_{j+1}\le \cos\Delta_j+(u/v)\sin\Delta_j.
\]
All right sides are positive. Their product is at least one. Write $\phi=\arctan(u/v)\in(0,\pi/2)$. The function
$\log(\cos\Delta+(u/v)\sin\Delta)
 =\log\sqrt{1+(u/v)^2}+\log\cos(\Delta-\phi)$
is strictly concave on the interval of positive factors in $(0,\pi)$. Jensen's inequality therefore gives
\[
 1\le\prod_j[\cos\Delta_j+(u/v)\sin\Delta_j]
 \le[\cos(2\pi/l)+(u/v)\sin(2\pi/l)]^l.
\]
As $l\ge3$, this implies $u/v\ge\tan(\pi/l)\ge\tan(\pi/M)$, proving \eqref{eq:r69-sector}.

At equality, $l=M$, every gap is $2\pi/M$, and all radius ratios equal one: the polygon is regular and every coordinate is one distinct vertex. The velocity at a vertex is along its outgoing edge. But
$(z-1)v_j=\sum_kQ_{jk}(v_k-v_j)$ is a nonnegative combination of chords. Strict convexity of this polygon makes all chords other than the self-chord and outgoing edge point strictly to the interior side of that edge. Their coefficients must therefore vanish. The remaining coefficient $t=Q_{j,j+1}$ satisfies $z-1=t(e^{2\pi i/M}-1)$ and is independent of $j$. Row normalization gives \eqref{eq:r69-lazy-cycle} and $t\le1$.
\end{proof}

\begin{theorem}[Exact first-sector memory threshold and optimizer]\label{thm:r69-sector-phase}
For $M\ge3$ and $0<\vartheta\le2\pi/M$, a response realization with at most $M$ live labels exists exactly when
\begin{equation}\label{eq:r69-sector-threshold}
 0<\alpha\le r_M(\vartheta)
 :=\frac1{\cos\vartheta+\sin\vartheta\tan(\pi/M)}.
\end{equation}
The threshold is attained. At $\alpha=r_M(\vartheta)$ the least memory is exactly $M$, the transition matrix is the lazy directed cycle in \eqref{eq:r69-lazy-cycle}, and, in cyclic coordinates with prescribed initial label zero, its decoder is uniquely
\[
 d_j=\tfrac12[1+\gamma\cos(\varphi+2\pi j/M)].
\]
This is global over arbitrary stochastic transition and decoder rows, not an optimization restricted in advance to circulants.
\end{theorem}
\begin{proof}
The spectral criterion and \eqref{eq:r69-sector} prove necessity. For sufficiency put $\omega=e^{2\pi i/M}$ and $J=\mathbf1\mathbf1^{\mathsf T}/M$. The nonnegative coefficients
\[
 b=\frac{\alpha\sin\vartheta}{\sin(2\pi/M)},\qquad
 a=\alpha\cos\vartheta-b\cos(2\pi/M),\qquad
 c=1-a-b
\]
sum to one under \eqref{eq:r69-sector-threshold}. The inequalities $a,b\ge0$ follow from $0<\vartheta\le2\pi/M$; $c\ge0$ is exactly the stated threshold. Then $Q=aI+bC_M+cJ$ has eigenvector $(1,\omega,\ldots,\omega^{M-1})$ at eigenvalue $\alpha e^{i\vartheta}$ and the displayed decoder realizes the response.

At the threshold, strict monotonicity of $\cot(\pi/M)$ excludes every smaller size. Lemma~\ref{lem:r69-sector} fixes $Q$. Here $b=t>0$, so $C_M=(Q-(1-t)I)/t$. The rows $e_0^{\mathsf T}Q^n$, $0\le n<M$, span all coordinate rows, because the rows $e_0^{\mathsf T}C_M^j$ do and each is a polynomial combination of the former. Thus all response coefficients determine the decoder uniquely. This proves the rigidity assertion including the endpoint $t=1$.
\end{proof}

For the inherited algebraic angle $\cos\vartheta=3/5$, $\sin\vartheta=4/5$, let
\[
 a_j=\frac5{3+4\tan(\pi/j)},\qquad j=3,4,5,6.
\]
One has $0<\vartheta<2\pi/6$. Consequently exact memory is three on $(0,a_3]$, four on $(a_3,a_4]$, five on $(a_4,a_5]$, and six on $(a_5,a_6]$. In particular $a_4=5/7$. Above $a_6$ at least seven labels are needed; no assertion that a regular polygon is always globally optimal beyond the first sector is made. The old three-label threshold is recovered rather than reproved as a new independent result.

\subsection{Exact worst-phase size and arithmetic improvements}
For an integer $q\ge3$ define the distance to the nearest $q$th root angle by
\[
 \delta_q(\vartheta)=\min_{j\in\mathbb Z}|\vartheta-2\pi j/q|
             \quad\text{modulo }2\pi,\qquad 0\le\delta_q\le\pi/q,
\]
and define the radial boundary of the regular polygon by
\begin{equation}\label{eq:r69-polygon-radius}
 R_q(\vartheta)=\frac{\cos(\pi/q)}{\cos(\pi/q-\delta_q(\vartheta))}.
\end{equation}

\begin{lemma}[Directional polygon construction and power obstruction]\label{lem:r69-arithmetic-tools}
If $0<\alpha\le R_q(\vartheta)$, a $q$-label exact response realization exists. Conversely, any $M$-label realization satisfies, for every positive integer $k$,
\begin{equation}\label{eq:r69-power-obstruction}
 \alpha^k\bigl[\cos(k\vartheta)+|\sin(k\vartheta)|\tan(\pi/M)\bigr]\le1.
\end{equation}
For $h=\pi/q$ and $0\le\delta=\delta_q\le h$,
\begin{equation}\label{eq:r69-polygon-error}
 1-R_q(\vartheta)\le2h\delta.
\end{equation}
\end{lemma}
\begin{proof}
The intersection of the ray at angle $\vartheta$ with the chord joining its two adjacent unit roots has radius \eqref{eq:r69-polygon-radius}. Thus $\alpha e^{i\vartheta}$ is in their triangle with zero. Write it as $a\omega^j+b\omega^{j+1}+c0$, with $a,b,c\ge0$ and sum one. The stochastic matrix $aC_q^j+bC_q^{j+1}+cJ$ has that eigenvalue, giving the first assertion by Theorem~\ref{thm:r69-spectral}. The matrix $Q^k$ of any realization is still stochastic and has eigenvalue $\alpha^ke^{ik\vartheta}$. Apply the sector lemma; for a real eigenvalue the displayed inequality is automatic. Finally $\cos(h-\delta)\ge\cos h\ge1/2$ and
$\cos(h-\delta)-\cos h=2\sin(h-\delta/2)\sin(\delta/2)\le h\delta$.
Division proves \eqref{eq:r69-polygon-error}.
\end{proof}

\begin{theorem}[Exact worst-phase memory]\label{thm:r69-worst-phase}
For every $0<\alpha<1$,
\begin{equation}\label{eq:r69-worst-phase}
 \max_{\sin\vartheta\ne0}m(\alpha,\vartheta)
  =\max\left\{3,\left\lceil\frac\pi{\arccos\alpha}\right\rceil\right\}.
\end{equation}
In particular the worst-phase exact label count is asymptotic to
$\pi/\sqrt{2(1-\alpha)}$ as $\alpha\uparrow1$. The phase may depend on $\alpha$; this is not a lower rate for each fixed phase.
\end{theorem}
\begin{proof}
Let the displayed integer be $K$. The regular $K$-gon contains the disk of radius $\cos(\pi/K)\ge\alpha$, so the polygon lemma gives the upper simultaneously for all phases. If $K=3$, three distinct nonzero modes force at least three states. Otherwise put $M=K-1\ge3$ and choose $\vartheta=\pi/M$. The definition of $K$ gives $\alpha>\cos(\pi/M)$; the exact first-sector threshold at this angle is $\cos(\pi/M)$. No $M$-label realization exists, proving the lower and attainment of the maximum. The asymptotic follows from $1-\cos t\sim t^2/2$.
\end{proof}

A number $x$ is called badly approximable here with constant $c>0$ when
\begin{equation}\label{eq:r69-diophantine}
 \|kx\|_{\mathbb R/\mathbb Z}\ge c/k\qquad(k\ge1).
\end{equation}
The constant is part of the hypothesis, not something inferred from a finite list of continued-fraction coefficients. The next theorem holds for every such phase and proves a different, matching fixed-phase rate.

\begin{theorem}[Sharp cubic-root memory at a fixed Diophantine phase]\label{thm:r69-cubic}
Let $x=\vartheta/(2\pi)$ satisfy \eqref{eq:r69-diophantine}. Then, for $0<\alpha<1$ and $\epsilon=1-\alpha$,
\begin{equation}\label{eq:r69-cubic-bounds}
 \left(\frac{\pi^2c}{16\epsilon}\right)^{1/3}
 \le m(\alpha,\vartheta)
 \le \max\left\{4,\left\lceil\frac3c\right\rceil,
              \left\lceil\left(\frac{4\pi^2}{c^2\epsilon}\right)^{1/3}\right\rceil\right\}.
\end{equation}
Thus $m(\alpha,\vartheta)=\Theta_c((1-\alpha)^{-1/3})$. The exact live-state bits are $\frac13\log_2(1/(1-\alpha))+O_c(1)$ in the ideal finite-label model. These are exact stochastic realizations at the original least expected-cost budget, not approximation of the response or additional physical observations. The statement does not supply matching program-bit or exact-sampling complexity bounds.
\end{theorem}
\begin{proof}
We use only the following elementary approximation fact. For any irrational $x$ and integer $Q\ge1$, among the $Q+1$ fractional parts $0,x,\ldots,Qx$ two fall in the same one of $Q$ equal intervals. Subtracting them gives $1\le k\le Q$ with $0<\|kx\|\le1/Q$.

For the lower, take a realizing size $M\ge3$ and apply that fact with $Q=2M$. Set $\delta=2\pi\|kx\|$ and $h=\pi/M$. Then
\[
 \pi c/M\le\delta\le h\le\pi/3.
\]
Writing $F=\cos\delta+\sin\delta\tan h$, we have $F\le\sec h\le2$. Also $\sin\delta\ge3\delta/4$ on $[0,\pi/3]$, $\tan h\ge h$ and $1-\cos\delta\le\delta^2/2$, so
\[
 F-1\ge3h\delta/4-\delta^2/2\ge h\delta/4.
\]
By \eqref{eq:r69-power-obstruction}, $\alpha^kF\le1$. Hence
\[
 2M\epsilon\ge k\epsilon\ge1-\alpha^k
       \ge1-1/F\ge h\delta/8
       \ge\pi^2c/(8M^2).
\]
This proves the lower constant.

For the upper, let $Q$ be the integer on the right of \eqref{eq:r69-cubic-bounds}. The approximation fact provides integers $1\le q\le Q$ and $a$ with $|qx-a|\le1/Q$. The lower bound \eqref{eq:r69-diophantine} forces $q\ge cQ\ge3$. The distance from $\vartheta$ to a $q$th-root angle is at most $2\pi/(qQ)$. The polygon-error bound gives
\[
 1-R_q(\vartheta)\le\frac{4\pi^2}{q^2Q}
       \le\frac{4\pi^2}{c^2Q^3}\le\epsilon.
\]
Thus $\alpha\le R_q(\vartheta)$ and the explicit circulant construction realizes the response with $q\le Q$ labels. Both bounds apply for the same fixed phase as $\epsilon\downarrow0$. Taking base-two logarithms and rounding the label bits proves the last assertion.
\end{proof}

The hypothesis is nonempty and has elementary physical examples. For
$x=(3-\sqrt5)/2$, choose the closest integer $a$ to $kx$. The integer
$a^2-3ak+k^2$ is nonzero. Since the other root is $x'=(3+\sqrt5)/2$ and
$|a/k-x'|\le\sqrt5+1/2<3$, factorization gives
$|a/k-x|>1/(3k^2)$. Thus \eqref{eq:r69-diophantine} holds with $c=1/3$. This provides a fixed qubit rotation with the matched exponent. Its trigonometric primitive entries are not asserted to be algebraic. In particular the Diophantine condition has not been proved here for the different algebraic rotation $(3+4i)/5$.

\begin{corollary}[A stronger fixed-algebraic-angle exclusion]\label{cor:r69-algebraic-angle}
For the inherited angle $e^{i\vartheta}=(3+4i)/5$ and $0<\alpha<1$, every exact $M$-label stochastic realization satisfies
\begin{equation}\label{eq:r69-algebraic-bound}
 1-\alpha\ge\frac3{16M^2\,25^M}.
\end{equation}
Consequently its exact memory is at least of order $\log(1/(1-\alpha))$ as $\alpha\uparrow1$. This improves the earlier $\log/\log\log$ lower, without assigning the cubic-root law to this angle or claiming a matching upper for it.
\end{corollary}
\begin{proof}
Again choose $1\le k\le2M$ with $\delta=2\pi\|k\vartheta/(2\pi)\|\le\pi/M$. The Gaussian integer $(3+4i)^k-5^k$ is nonzero: otherwise the unit complex number would be a root of unity, forcing $6/5$ to be a rational algebraic integer. Its modulus is at least one. Therefore
$\delta\ge|1-e^{ik\vartheta}|\ge5^{-k}\ge25^{-M}$.
Repeat the lower proof of Theorem~\ref{thm:r69-cubic} up to $2M(1-\alpha)\ge h\delta/8$, with $h=\pi/M$. This gives the stronger constant $\pi/(16M^2 25^M)$ and hence \eqref{eq:r69-algebraic-bound}. Taking logarithms gives
$\log(1/(1-\alpha))\le\log(16/3)+2\log M+M\log25$, proving the order assertion.
\end{proof}

\subsection{Finite spectral mixtures and higher-dimensional protected systems}
The matching exponent is not confined to a two-dimensional physical carrier. Let $r\ge1$, let $\vartheta_1,\ldots,\vartheta_r\in(0,\pi)$ be distinct, and let $w_j>0$ with $\sum_jw_j\le1$. For fixed phases $\varphi_j$, consider
\begin{equation}\label{eq:r69-mixture}
 b_n^{\rm mix}=\tfrac12\left[1+\alpha^n
             \sum_{j=1}^r w_j\cos(n\vartheta_j+\varphi_j)\right].
\end{equation}
Write $m_{\rm mix}(\alpha)$ for its least pure-start stochastic memory.

\begin{theorem}[Spectral mixtures preserve the matched arithmetic exponent]\label{thm:r69-finite-mixture}
For $0<\alpha<1$, the exact memory of \eqref{eq:r69-mixture} satisfies
\begin{equation}\label{eq:r69-mixture-bounds}
 \max\{2r+1,\max_j m(\alpha,\vartheta_j)\}
 \le m_{\rm mix}(\alpha)
 \le 1+\one_{\{\sum_jw_j<1\}}+\sum_{j=1}^r m(\alpha,\vartheta_j).
\end{equation}
All models are realized with the counted prescribed initial label. If every fixed $\vartheta_j/(2\pi)$ has a global badly-approximable certificate, then
$m_{\rm mix}(\alpha)=\Theta((1-\alpha)^{-1/3})$, with constants depending on the finite phase list. The same conclusion holds when some phases are rational and at least one is certified badly approximable. This probability response is physically prepared on dimension at most $2r+1$, and its exact memory is the boundary-optimal observer memory of the same squared task after the standard two-dimensional herald-flag embedding.
\end{theorem}
\begin{proof}
The $2r+1$ modes $1,\alpha e^{\pm i\vartheta_j}$ are pairwise distinct and have nonzero coefficients. Their Hankel rank is $2r+1$ by Vandermonde independence. Every realizing characteristic polynomial vanishes at every nonconstant mode, giving both lower bounds through Theorem~\ref{thm:r69-spectral}.

For each component choose a minimal stochastic realization $(Q_j,d_j)$ starting in its first label. Keep their state sets disjoint and add one initial label $o$ with decoder $b_0^{\rm mix}$. On its first failure, its row enters block $j$ with subprobability row $w_je_1^{\mathsf T}Q_j$. If $w_0=1-\sum_jw_j>0$, its remaining mass enters one absorbing label with decoder $1/2$. Within each block use $Q_j$. Thus after $n\ge1$ failures the output probability is
$w_0/2+\sum_jw_je_1^{\mathsf T}Q_j^nd_j=b_n^{\rm mix}$, while the initial decoder handles $n=0$. No physical call is added and no persistent choice is stored outside the counted labels. This proves the upper. Theorem~\ref{thm:r69-cubic} supplies matched exponents for the finite certified phase list. For a rational phase of order $q$, the matrix $\alpha C_q+(1-\alpha)J$ gives a uniform $q$-label upper, so adding those components does not change the exponent.

For the physical realization, take a direct sum of $r$ qubits, prepare the phase-rotated pure state in block $j$ with weight $w_j$, and use the damped rotation channel in that block. Drop off-diagonal block coherences. If $w_0>0$, add one scalar block of mass $w_0$ with effect $1/2$ and identity evolution. Each diagonal branch is completely positive and their direct sum is trace preserving. The terminal effect is the direct sum of the qubit $Z$ effects and the scalar half effect. Its probability is exactly \eqref{eq:r69-mixture}. Theorem~\ref{thm:r67-embedding} therefore gives the same fixed squared loss and cost $1/p$, and forces this response at the optimum. The observer's independently randomized block selection does not reveal or copy the protected block.
\end{proof}

The three memory statements have different quantifiers: the exact worst-phase law permits the phase to vary with the noise, the cubic-root law fixes any phase with the displayed arithmetic certificate, and the last corollary fixes the older algebraic rotation without such a certificate. Noise-dependent classical memory bounds for quantum processes were already established by Fanizza, Lumbreras and Winter \cite{Fanizza24}; their worst-phase process bound is not a fixed-phase matched law for this response task. The stochastic regions and their Type-0 rigidity are also classical. The present proofs make the task-level reduction and the new arithmetic comparison explicit, rather than claiming earliest priority from a literature search. No statement in this section solves the different squared-sign prediction frontier or the phase-type timing frontier.

\section{Finite realization of uniformly interior instruments}\label{sec:interior-instruments}
The boundary decision theorems identify the observer memory needed for an optimal task. We now consider a distinct, stronger form of physical-model reduction: one finite hidden classical model must preserve every controlled report word, every stopping rule and every terminal measurement of a declared instrument family. Its hidden labels are not made available to the observer. The distinction is essential: a finite hidden-state realization does not in general make the posterior, or an optimal observer, finite-state.

Invariant positive cones are classical realization tools \cite{Heller65,MonrasWinter14}. Our purpose here is not another abstract cone criterion. We derive a finite invariant set from a quantitative condition on every actual report branch, give its dimension dependence and an explicit sparse implementation, and prove exact stopped-law and cost transfer. The commands need not commute, need not reset the input, and need not share a stationary state. The condition is sufficient, not necessary for every finite hidden realization. A matching noise exponent below concerns the expressly defined state-positive class, not every possible classical response representation.

\subsection{A finite hull containing all uniformly mixed states}
Let $\mathcal D_d$ be the density matrices on $\mathbb C^d$, $d\ge2$, and put
\[
 \mathcal D_d(\eta)=\{\rho\in\mathcal D_d:\rho\succeq\eta I/d\}
       =(1-\eta)\mathcal D_d+\eta I/d,\qquad 0<\eta\le1.
\]
For unit vectors modulo phase use the metric
\[
 \operatorname{dist}([u],[v])
 =\sqrt{1-|\langle u,v\rangle|^2}
 =\tfrac12\|uu^*-vv^*\|_1=\|uu^*-vv^*\|_{\rm op}.
\]
In particular it satisfies the triangle inequality. All convex hulls below are hulls of actual positive density matrices, not signed coordinates.

\begin{lemma}[An explicit interior hull]\label{lem:r68-hull}
There are pure states $P_1,\ldots,P_K$ such that
\begin{equation}\label{eq:r68-hull}
 \mathcal D_d(\eta)\subseteq\operatorname{conv}\{P_1,\ldots,P_K\}
              \subseteq\mathcal D_d,
 \qquad K\le\left\lceil(4d/\eta)^{d-1}\right\rceil.
\end{equation}
Any additional $k$ prescribed initial states may be included as $k$ further generators. The generators are shared by all maps satisfying the same interior margin.
\end{lemma}
\begin{proof}
For the unitarily invariant probability measure on complex projective space, a ball of radius $r\in[0,1]$ has mass $r^{2d-2}$. To verify the formula, normalize $d$ independent standard complex Gaussian coordinates. Their squared moduli before normalization are independent exponential variables. Changing variables to their sum and normalized proportions gives constant density on the probability simplex. Thus $|u_1|^2$ has density $(d-1)(1-x)^{d-2}$ on $[0,1]$, and the probability of $|u_1|^2>1-r^2$ is $r^{2d-2}$. This also proves the formula at the endpoints by continuity.

Set $\varepsilon=\sqrt{\eta/d}$. A maximal set of rays with pairwise distances greater than $\varepsilon$ is an $\varepsilon$-net. The radius-$\varepsilon/2$ open balls about its points are disjoint. Their equal positive masses bound the cardinality by $(4/\varepsilon^2)^{d-1}$. Such a finite maximal set exists by successively adding uncovered points; the same mass bound forces this process to stop.

Let $A$ be any Hermitian matrix, with greatest and least eigenvalues $a,b$. Choose a net vector $v$ within $\varepsilon$ of a greatest-eigenvalue vector. Expanding in that eigenbasis gives
\[
 \tr(Avv^*)\ge a-(a-b)\varepsilon^2.
\]
On the other hand,
\[
 \max_{\rho\in\mathcal D_d(\eta)}\tr(A\rho)
 =a-\eta\bigl(a-\tr A/d\bigr)
 \le a-\eta(a-b)/d.
\]
The net hull therefore has support function at least that of $\mathcal D_d(\eta)$ in every Hermitian direction. Finite-dimensional separation proves containment. The other inclusion and the additional-generator assertion follow from positivity and convexity.
\end{proof}

\subsection{Report-resolved positive interpolation}
There is a finite set of globally legal commands $\mathcal A$ and, for each command, a finite report set $Y_a$. Write $\mathcal I_{a,y}$ for its completely positive report map and assume $\sum_y\mathcal I_{a,y}$ is trace preserving. A finite list $\rho^{(1)},\ldots,\rho^{(k)}$ of allowed prepared states is specified. The initial preparation index is fixed before execution; it is not an unknown state revealed to the observer. Assume the following unnormalized inequality for every command, report and positive input:
\begin{equation}\label{eq:r68-branch-floor}
 \mathcal I_{a,y}(X)\succeq
       \frac{\eta}{d}\tr\bigl(\mathcal I_{a,y}(X)\bigr)I,
 \qquad X\succeq0.
\end{equation}
A zero branch is allowed. No probability is divided out at zero. The hypothesis is on each recorded branch, not merely on the sum of the maps. It is satisfied when a depolarizing channel with mixing weight at least $\eta$ is physically applied after each report branch. Merely assuming positive definite outputs at each nonzero input is not used to infer a uniform ratio at a vanishing branch.

A \emph{state-positive realization} consists of density matrices $\tau_1,\ldots,\tau_M$, including the prescribed preparations as designated initial labels, and nonnegative matrices $Q_{a,y}$ satisfying
\begin{align}
 \mathcal I_{a,y}(\tau_i)&=\sum_jQ_{a,y}(i,j)\tau_j,
       &\sum_{y,j}Q_{a,y}(i,j)&=1.\label{eq:r68-intertwining}
\end{align}
The second equation includes all reports of a fixed command. The simulator in state $i$, on receiving command $a$, draws the pair $(y,j)$ with these probabilities, reports only $y$, and retains $j$ privately. For a terminal positive-operator-valued measurement $\{E_z\}_z$, it reports $z$ with probability $\tr(E_z\tau_j)$. A fixed finite decoder dictionary is handled by the unchanged external observer.

\begin{theorem}[Uniform finite controlled realization]\label{thm:r68-instrument}
Under \eqref{eq:r68-branch-floor}, a state-positive realization exists with
\begin{equation}\label{eq:r68-size}
 M\le k+\left\lceil(4d/\eta)^{d-1}\right\rceil.
\end{equation}
Each row of each $Q_{a,y}$ may be chosen with at most $d^2$ nonzero entries. The count is independent of the horizon, number of commands, number of reports, and choice of terminal measurement; those quantities still affect the transition-table size.

For every adaptive randomized command policy using only past declared reports and its own randomness, this one realization preserves all finite report laws and all adaptively chosen terminal measurement laws. It also preserves the complete stopped law, including the event of never stopping. Every nonnegative cost that is a function of the executed command/report path has the same expectation, finite or infinite. If calls instead have specified unobserved nonnegative mean observables $G_a$, the assignment $c_a(i)=\tr(G_a\tau_i)$ preserves their total expected cost at every stopping rule. This latter assertion is an expectation statement, not a joint-law assertion about an unspecified random cost observation.
\end{theorem}
\begin{proof}
Take the hull of Lemma~\ref{lem:r68-hull} and append the prepared states. For any generator $\tau_i$, let $s=\tr\mathcal I_{a,y}(\tau_i)$. If $s=0$, positivity makes the branch matrix zero, and set that whole row of $Q_{a,y}$ to zero. If $s>0$, its normalized output belongs to $\mathcal D_d(\eta)$ and hence to the pure-state core hull. Choose convex coefficients there and multiply them by $s$. This proves the first identity. Taking its trace and summing in $y$ proves row normalization.

A convex combination in the trace-one Hermitian affine space needs at most $d^2$ generators. Indeed any larger supported list is affinely dependent; subtract a scalar multiple of an affine dependence until one positive coefficient becomes zero, preserving nonnegativity, total mass and represented matrix, and iterate. Applying this finite reduction separately to each branch gives the sparsity claim.

For a fixed command/report history $h$, let $X_h$ be its unnormalized physical state and $p_h$ the simulator's unnormalized hidden-state row. Induction using \eqref{eq:r68-intertwining} gives the matrix identity
\begin{equation}\label{eq:r68-history}
 X_h=\sum_i p_h(i)\tau_i.
\end{equation}
It holds initially because the initial label is the prescribed preparation. Multiplying both sides by a report branch gives the next identity. Traces give the report probability, and pairing with any terminal effect gives its joint audit probability.

For an adaptive observer the command probabilities and commitment probabilities conditional on a declared history are the same in the two experiments. Condition on its private tape, apply the same identity at each finite history, and integrate over that tape. Equivalently keep one unnormalized identity for every observer-memory state. The simulator's hidden state is not an extra input to this policy. Thus all finite stopped cylinders agree. The events of survival through $n$ calls decrease to nontermination; their equal probabilities give equality there as well. Bounded terminal measures follow by summing disjoint finite stopping events. Nonnegative path costs follow from equality of the path law and monotone convergence.

For mean observables, the expected cost contribution at a begun history is $\tr(G_aX_h)=\sum_i p_h(i)c_a(i)$. Sum these identities first to a finite horizon and then monotonically to the stopping time. No total-variation bound on bounded terminal scores is used to infer an unbounded waiting-cost bound.
\end{proof}

This theorem realizes the prepared experiment, not an arbitrary entangled quantum input channel. It gives no affine encoding of every density matrix into one simplex and does not grant the hidden simulator label to the original observer. For a single prescribed initial state one anchor suffices. If $d=1$, the original scalar instrument already has one state. If $\eta=1$, every nonzero output is $I/d$, so one core state plus the preparation anchors suffices without a projective net.

\begin{corollary}[Same-policy resource and risk invariance]\label{cor:r68-risk}
The original instrument and its hidden classical realization have identical feasible expected-cost budgets and identical risks for each common report-based observer. Consequently their optimal risk curves agree at every fixed observer memory bound and every fixed decoder dictionary, including feasible boundary budgets. The assertion also holds modelwise for an unknown family satisfying one uniform margin, with model-dependent hidden transition tables and the same report interface; the parameter is not supplied to the observer. If the prepared state varies with the model, it is represented by one designated model-dependent anchor and a common pure-state core.

An observer with $\ell$ live labels coupled to an $M$-state simulator has at most $\ell M$ joint live states. These are joint system states, not a claim that $M$ is the observer's optimal memory. Giving the hidden label to the observer changes the statistical experiment and invalidates the claimed same-policy comparison.
\end{corollary}
\begin{proof}
Theorem~\ref{thm:r68-instrument} equates the cost and risk of every admissible observer without changing its transition rows, dictionary or report information. Taking the same constrained infimum gives the assertion, even when the infimum is not attained in that memory class. For a model family use the same core hull in every model and solve its branch identities separately; each modelwise equality holds for the same observer. Taking the supremum over the unknown model therefore preserves the robust objective and constraints. The product count records the two distinct retained coordinates.
\end{proof}

\subsection{A matching exponent for state-positive universal realization}
Let $v_d(\eta)$ be the least number of arbitrary density matrices whose convex hull contains $\mathcal D_d(\eta)$. Mixed vertices are permitted. This number measures uniform state-positive representation, not the smallest hidden model for one selected scalar response.

\begin{theorem}[Sharp noise exponent at fixed physical dimension]\label{thm:r68-noise-exponent}
For $d\ge2$ and $0<\eta\le1/2$, put $s=\eta(1-1/d)$. Then
\begin{equation}\label{eq:r68-noise-count}
 s^{-(d-1)}\le v_d(\eta)
           \le\left\lceil(4d/\eta)^{d-1}\right\rceil.
\end{equation}
Thus the exponent $d-1$ as $\eta\downarrow0$ is necessary and sufficient at each fixed $d$. Constants in this comparison depend on $d$.

For every $0<\eta\le1/4$ there is also a finite-command, unit-call-cost prepared instrument satisfying \eqref{eq:r68-branch-floor} for which every state-positive exact realization preserving all terminal effects needs at least
\begin{equation}\label{eq:r68-finite-noise-lower}
 \bigl(4/(5\eta)\bigr)^{d-1}
\end{equation}
hidden states. Its command count may grow as the margin decreases. No lower bound of this exponent for unrestricted, possibly non-state-positive, response realizations is asserted.
\end{theorem}
\begin{proof}
The upper is Lemma~\ref{lem:r68-hull}. For the lower, for each unit $u$ consider the state $(1-\eta)uu^*+\eta I/d$. Pairing a convex representation in density vertices $\tau_j$ with $uu^*$ shows that at least one vertex satisfies $\langle u,\tau_ju\rangle\ge1-s$.

For any fixed density matrix $\tau$, this inequality holds on a set of projective measure at most $s^{d-1}$. If its largest eigenvalue $a$ is below $1-s$, the set is empty. Otherwise $a\ge1-s>1/2$ and its top eigenvector $v$ is unique. With $x=|\langle u,v\rangle|^2$, all other eigenvalues are at most $1-a$, so
\[
 \langle u,\tau u\rangle\le ax+(1-a)(1-x).
\]
To reach $1-s$ one needs $x\ge(a-s)/(2a-1)\ge1-s$; the last inequality is equivalent to $(1-a)(1-2s)\ge0$. The projective ball formula from the lemma now bounds the measure by $s^{d-1}$. The caps associated with all vertices cover every ray, so their measures sum to at least one. This proves the first lower.

For the finite-command assertion take a finite projective $\eta/4$-net of rays $u_j$ and a command for each replacement channel
\[
 \Phi_j(X)=\tr(X)\bigl((1-\eta)u_ju_j^*+\eta I/d\bigr).
\]
These are valid completely positive, trace-preserving maps with the required margin, one deterministic report and unit cost. Exactness of every terminal effect after that command forces its output matrix to be a convex combination of the simulator's density vertices. Thus at each net point some vertex has expectation at least $1-s$. For an arbitrary ray $u$ choose a net point within $\eta/4$. Since the operator norm distance of the two projectors is at most $\eta/4$, pairing with any density matrix changes the expectation by at most $\eta/4$. The caps at level $1-s-\eta/4$ therefore cover all rays. Here $s+\eta/4\le5\eta/4<1/2$, and the preceding cap bound gives at least $(s+\eta/4)^{-(d-1)}\ge(4/(5\eta))^{d-1}$ vertices. This is a finite physical family, not an uncountable list disguised as a finite input.
\end{proof}

For unrestricted classical hidden models a single scalar response may admit non-quantum intermediate response vectors. The state-positive lower does not exclude that possibility. Conversely, a hidden decoder defined as a positive normalized affine functional of every terminal effect is necessarily of the form $E\mapsto\tr(E\tau)$ for a density matrix $\tau$: extend the affine functional linearly from the effect interval to the Hermitian space, use the trace pairing, and test rank-one positive effects. This explains the positivity requirement in the count rather than hiding it in an abstract factorization.

\subsection{Sparse algebraic construction and several-step interiority}
\begin{proposition}[Finite algebraic construction with sparse branch rows]\label{prop:r68-algebraic}
For a finite instrument family and finite preparation list with real-algebraic entries and a supplied algebraic margin $\eta>0$, a state-positive realization can be constructed by a terminating real-algebraic procedure. One may use the size bound \eqref{eq:r68-size}, and all generator matrices and transition probabilities can be chosen algebraic. At most $d^2$ successors per command/report row are needed. The original margin condition is itself a finite quantified real-algebraic condition on pure inputs; a positive rational margin, if one exists, can be found by enumeration and such tests.

This is a finite synthesis assertion, not a polynomial-time bound. The retained-state count, the number and algebraic encoding lengths of table entries, and the resources of exact random sampling remain different quantities.
\end{proposition}
\begin{proof}
The existence proof uses at most $K=\lceil(4d/\eta)^{d-1}\rceil$ pure core projectors and the stated anchors. Treat their real Hermitian coordinates and all branch coefficients as variables. Impose $P_j\succeq0$, $\tr P_j=1$, $P_j^2=P_j$, the fixed anchor identities, nonnegative normalized rows, and the finitely many matrix equations \eqref{eq:r68-intertwining}. These are polynomial conditions over the real algebraic field; positivity can be expressed by principal minors. The preceding theorem proves feasibility. Real-closed-field elimination and sample-point construction give an algebraic solution \cite{Seidenberg54}. Row sparsity can be obtained afterwards by elimination of an affine dependence over the same algebraic field, or by a finite disjunction of supports of size at most $d^2$.

For margin testing, it suffices by linearity to quantify over pure inputs $uu^*$ with $\|u\|=1$; the output inequality is again a principal-minor condition in real and imaginary coordinates. If some real positive margin works, a smaller positive rational margin works. Enumeration of positive rational candidates with an exact test therefore eventually succeeds on the strict-interiority class. No termination is asserted when no positive margin exists. For $J=\sum_a|Y_a|$, a sparse implementation stores at most $JM d^2$ nonzero branch probabilities, besides $Md^2$ density coordinates and any chosen audit table. These are storage counts of mathematical entries, not bit bounds for their encodings.
\end{proof}

\begin{theorem}[Interiority after a finite block, with its clock counted]\label{thm:r68-block}
Suppose instead that for some integer $L\ge1$, every length-$L$ command/report branch product satisfies \eqref{eq:r68-branch-floor} with the same margin. Individual branches need not do so. Put $J=\sum_a|Y_a|$. Then an exact state-positive realization exists with at most
\begin{equation}\label{eq:r68-block-count}
 \left(k+\left\lceil(4d/\eta)^{d-1}\right\rceil\right)
                        \sum_{j=0}^{L-1}J^j
\end{equation}
hidden states. It preserves intermediate reports and stopping within the block; it does not replace $L$ calls by one unpriced call.
\end{theorem}
\begin{proof}
Start with the core projectors and preparation anchors. For each, list its normalized nonzero images under all command/report products of length less than $L$, retaining the product as part of the hidden label. There are at most the displayed number. At a depth below $L-1$, an actual branch goes to its listed child with its actual conditional probability. At depth $L-1$, the next product output is either zero or belongs to $\mathcal D_d(\eta)$; decompose it into the core and restart at depth zero. Zero branches have zero rows. These transitions satisfy the same matrix interpolation identity at every single call. Theorem~\ref{thm:r68-instrument}'s history induction therefore applies without suppressing any intermediate observation, audit or charge. The depth is encoded in the enlarged state label, not in an external clock.
\end{proof}

\subsection{Why forgetting the report invalidates the premise}
The average channel can be completely depolarizing even though every recorded branch sends pure states to pure states. The Pauli instrument $\mathcal I_P(X)=PXP^*/4$ on a qubit is an example. Thus an average-state calculation cannot verify \eqref{eq:r68-branch-floor}.

A stronger example has, in addition, a recorded irrational rotation. Mix half of that Pauli instrument with the single recorded branch $\mathcal I_u(X)=UXU^*/2$, where
\[
 U=5^{-1/2}\begin{pmatrix}2&-1\\1&2\end{pmatrix}.
\]
Its average sends every density matrix to at least $I/4$, but the report word $u^n$ retains the pure state $U^n|0\rangle\langle0|U^{*n}$. These states are all distinct: $e^{i\vartheta}=(3+4i)/5$ is not a root of unity, since $2\cos\vartheta=6/5$ would otherwise be a rational algebraic integer. A finite state-positive simulator preserving all terminal effects would represent every one of these pure states as a convex combination of its finite density vertices. An extreme point of $\mathcal D_2$ can only be represented by identical density matrices, so every distinct orbit state would have to occur as a vertex. This is impossible. Thus average interiority does not even guarantee finite state-positive realization, while all finite report words here have well-defined nonzero probabilities.

The example is a boundary qualification of the positive theorem. It is not used in place of that theorem, and no unrestricted classical hidden-state impossibility is inferred from this particular extreme-state proof alone.

\subsection{An explicit unrestricted memory obstruction at weak damping}
The preceding matching exponent is for state-positive realization. A different argument gives a quantitative, albeit non-sharp, obstruction for every stochastic realization of the fixed-angle scalar qubit response. It sharpens the qualitative divergence in Theorem~\ref{thm:qubit-positive-phase} without changing its exact three-label result.

\begin{lemma}[A finite near-peripheral spectral witness]\label{lem:r68-cycle}
If an $M\times M$ row-stochastic matrix has a nonzero eigenvalue $z=\alpha e^{i\theta}$, $0<\alpha\le1$, then either $1-\alpha>1/(2M)$ or some $1\le l\le M$ satisfies
\begin{equation}\label{eq:r68-cycle}
 |1-e^{il\theta}|\le
       \bigl(4\sqrt{2M^{M+1}}+M\bigr)\sqrt{1-\alpha}.
\end{equation}
The assertion does not require irreducibility, reversible rows or a lower bound on all positive transition probabilities.
\end{lemma}
\begin{proof}
Normalize an eigenvector $x$ so that its largest coordinate modulus is one, attained at $i_0$, and multiply its phase so $x_{i_0}=1$. Let $X_n$ be the Markov chain starting there. Then $\mathbb E x_{X_n}=z^n$ and
\[
 \mathbb E|x_{X_n}-z^n|^2
       =\mathbb E|x_{X_n}|^2-|z|^{2n}\le1-\alpha^{2n}.
\]
Choose once for each state a successor with transition probability at least $1/M$, and follow this directed path. Its first $n$ steps have probability at least $M^{-n}$, so their endpoint $i_n$ obeys
\[
 |x_{i_n}-z^n|\le\sqrt{M^n(1-\alpha^{2n})}
                \le\sqrt{2nM^n(1-\alpha)}.
\]
Two of $i_0,\ldots,i_M$ agree, say $i_j=i_k$ with $j<k$. Their two bounds give
$\alpha^j|1-z^{k-j}|\le2\sqrt{2M^{M+1}(1-\alpha)}$.
If $1-\alpha\le1/(2M)$, Bernoulli's inequality gives $\alpha^j\ge1/2$. Finally
$|1-e^{i(k-j)\theta}|\le|1-z^{k-j}|+1-\alpha^{k-j}$ and
$1-\alpha^{k-j}\le M(1-\alpha)\le M\sqrt{1-\alpha}$.
These inequalities prove the claim.
\end{proof}

\begin{corollary}[A rational exclusion margin for every memory size]\label{cor:r68-qubit-lower}
Fix $\zeta=(3+4i)/5$ and $0<\alpha<1$. If
\[
 \tfrac12\bigl[1+\alpha^n\operatorname{Re}(\zeta^n)\bigr]
           =e_1^{\mathsf T}Q^n d\quad(n\ge0),
 \qquad Q\mathbf1=\mathbf1,\quad Q\ge0,\quad 0\le d\le1,
\]
with $M$ classical states, then
\begin{equation}\label{eq:r68-qubit-lower}
 1-\alpha\ge\frac1{49M(25M)^M}.
\end{equation}
Consequently the exact memory in the fixed-angle calibration and common-squared-audit embeddings diverges at least on the order of
$\log(1/(1-\alpha))/\log\log(1/(1-\alpha))$ as $\alpha\uparrow1$. This lower is not claimed to match the retained polygon upper or to solve the general-$M$ frontier.
\end{corollary}
\begin{proof}
The response has three distinct nonzero exponential modes $1,\alpha\zeta,\alpha\bar\zeta$. Applying the characteristic recurrence of $Q$ to it and using Vandermonde independence forces $\alpha\zeta$ to be an eigenvalue. For every positive integer $l$, the nonzero Gaussian integer $(3+4i)^l-5^l$ has modulus at least one, so $|1-\zeta^l|\ge5^{-l}$. The nonzero assertion follows from the root-of-unity argument above.

If $1-\alpha\le1/(2M)$, Lemma~\ref{lem:r68-cycle} gives
\[
 5^{-M}\le\bigl(4\sqrt{2M^{M+1}}+M\bigr)\sqrt{1-\alpha}
             \le7M^{(M+1)/2}\sqrt{1-\alpha}.
\]
Squaring proves \eqref{eq:r68-qubit-lower}. In the other case that bound is weaker than $1/(2M)$ and is automatic. Taking logarithms gives the exact necessary inequality
\[
 \log\frac1{1-\alpha}\le\log49+\log M+M\log(25M),
\]
which implies the stated asymptotic order by elementary inversion.
\end{proof}

Noise-dependent classical-memory obstructions for quantum processes precede this result. In particular Fanizza, Lumbreras and Winter prove a worst-angle lower bound for noisy qubit processes, conditional on positive realizations for the whole phase family \cite{Fanizza24}. Their statement is not a lower rate for the one fixed angle in this corollary, nor does our weaker fixed-angle estimate replace it. The new general contribution of this section is the branch-resolved constructive realization and its stopped-resource identity, together with the state-positive dimension exponent. The established cone and spectral mechanisms are credited rather than renamed.

\section{Boundary stopping and finite stochastic realization}\label{sec:stochastic-boundary}
The weighted potential identity and the cost-active support criterion of Theorem~\ref{thm:weighted-rigidity} already allow model- and state-dependent reports. We now replace that sufficient coverage condition by a necessary-and-sufficient stopping criterion. We first characterize rigidity for prepared instruments with model- and state-dependent report probabilities. On the resulting boundary face, ties lead to a stochastic realization problem, not in general a deterministic residual quotient. The optimal-face and anchored convex realization theorems are inherited in full as Theorems~\ref{thm:weighted-face} and~\ref{thm:convex-memory}. Their pointed-polytope formulation below extends them to the larger rigidity class; neither the convex realization method nor the original robust-tie classification is claimed as new in this continuation.

\subsection{A primitive class and an exact boundary stopping criterion}
Let $\Theta$ be compact metric and let the protected physical space have finite dimension $d$. There is one globally legal command, a finite failure alphabet $Y$, and one productive herald $s$. Model $\theta$ has a prepared density matrix $\rho_\theta$ and completely positive maps $\mathcal I_{\theta,y}$ and $\mathcal I_{\theta,s}$ whose sum is trace preserving. All these data are continuous in $\theta$. A nonnegative call cost has conditional mean $\tr(G_\theta\rho)$, with continuous Hermitian $G_\theta$ satisfying
\[
 0<c_*I\preccurlyeq G_\theta\preccurlyeq c^*I<\infty.
\]
A failure forbids commitment. A herald permits commitment to a member of one fixed finite dictionary $C$, or rejection and continuation from the actual post-herald state. The full report is the only physical observation available to the controller. No extra cost observation, reset, state inspection or free clock is supplied. After commitment an unobserved prospective audit has continuous loss observables $0\preccurlyeq F_{\theta,r}\preccurlyeq LI$, $r\in C$. A separately executed fixed final-audit cost is added to all budgets.

Put $\mathcal T_\theta=\sum_y\mathcal I_{\theta,y}$ and impose the primitive transience test
\begin{equation}\label{eq:r67-transience}
 I-(\mathcal T_\theta^*)^dI\succ0\quad(\theta\in\Theta).
\end{equation}
This condition does not require scalar herald effects. Define
\[
 \delta=\min_\theta\lambda_{\min}\bigl(I-(\mathcal T_\theta^*)^dI\bigr)>0,
 \qquad \kappa=c^*d/\delta,
 \qquad H_\theta=\sum_{n\ge0}(\mathcal T_\theta^*)^nG_\theta.
\]
For $T\ge0$, set $\beta_T=(1-\delta)^{\lfloor T/d\rfloor}$, interpreting the zeroth power as one even when $\delta=1$. For a failure word $w=y_1\cdots y_n$, let $\mathcal I_{\theta,w}=\mathcal I_{\theta,y_n}\cdots\mathcal I_{\theta,y_1}$, with the identity map at the empty word. Its first-herald mass and terminal score are
\begin{align}
 \lambda_\theta(w)&=\tr\bigl(\mathcal I_{\theta,s}\mathcal I_{\theta,w}(\rho_\theta)\bigr),\label{eq:r67-word}\\
 g_\theta(w,r)&=\tr\bigl(F_{\theta,r}\mathcal I_{\theta,s}\mathcal I_{\theta,w}(\rho_\theta)\bigr).
\end{align}
In particular $0\le g_\theta(w,r)\le L\lambda_\theta(w)$. These are unnormalized quantities; null reports are never divided out.

\begin{lemma}[Uniform first-herald control from the instrument]\label{lem:r67-first-herald}
First-herald commitment is proper from every physical state in every model, has conditional remaining cost at most $\kappa$, and has survival mass after $T$ failures at most $\beta_T$. Its expected cost from the declared preparation is $t_\theta=\tr(H_\theta\rho_\theta)$, continuous in $\theta$. The least robust feasible cost is
\[
 N_0=\max_{\theta\in\Theta}t_\theta.
\]
Moreover, \eqref{eq:r67-transience} is equivalent, modelwise, to absence of a nonzero subspace invariant under every failure Kraus operator and annihilated by the herald effect.
\end{lemma}
\begin{proof}
For one model write $B=I-\mathcal T^*I=\mathcal I_s^*I$ and $O_n=I-(\mathcal T^*)^nI$. Then $O_{n+1}=B+\mathcal T^*O_n$. For Kraus operators $L_j$ of $\mathcal T$, positivity gives
$\ker O_{n+1}=\ker B\cap\bigcap_jL_j^{-1}(\ker O_n)$.
The decreasing kernels stabilize at their first equality and can strictly decrease at most $d$ times. Their stabilized space is exactly the failure-invariant dark subspace in the statement. This proves the equivalence without choosing a normalized posterior.

Continuity and compactness give $\delta>0$. The inequality $(\mathcal T_\theta^*)^dI\preccurlyeq(1-\delta)I$ iterates, uniformly over $\theta$, to the displayed survival bound. Summing in blocks of $d$ calls bounds expected cost by $c^*d/\delta$. The cost series converges uniformly in operator norm; hence $H_\theta$ and $t_\theta$ are continuous. Every legal committed trajectory contains all calls through its first herald. Nonnegative additional costs therefore make its expectation at least $t_\theta$ in that model. First-herald commitment attains all these modelwise costs simultaneously, proving the formula for $N_0$ and its attainability.
\end{proof}

Let $\mathcal W=\{w:\lambda_\theta(w)>0\text{ for some }\theta\}$. Define the normalized unused-budget margin
\begin{equation}\label{eq:r67-margin}
 a(w)=\inf_{\theta:\lambda_\theta(w)>0}
             \frac{N_0-t_\theta}{\lambda_\theta(w)}\quad(w\in\mathcal W).
\end{equation}
This is a criterion on the primitive first-herald experiment, not a condition on an unknown optimal decision policy.

\begin{theorem}[Necessary and sufficient boundary stopping rigidity]\label{thm:r67-rigidity}
Every robust policy of expected cost at most $N_0$ commits at the first herald in every model, up to null histories, if and only if $a(w)=0$ for every $w\in\mathcal W$.

More precisely, if $r_w$ is its conditional probability of rejecting its first herald after $w$, then
\begin{equation}\label{eq:r67-rejection-bound}
 c_*r_w\le a(w).
\end{equation}
If $a(w)>0$ for some word, there is a finite-label boundary-feasible policy that rejects that first herald with positive probability and then commits at the next herald. Thus failure of the criterion has an actual legal, paid implementation.
\end{theorem}
\begin{proof}
Before the first herald the sole command is fixed. Conditional on its full word, the private controller randomness is independent of the physical state and of the model: it changes no executed instrument. Consequently the rejection probability $r_w$ is the same in all models. On that event at least one further call is required, whose conditional expected cost is at least $c_*$. Decomposing the nonnegative cost at the first herald gives
\[
 \E_\theta C_\tau\ge t_\theta+
                 c_*\sum_w\lambda_\theta(w)r_w.
\]
Keeping one word and using feasibility proves \eqref{eq:r67-rejection-bound}. If every $a(w)$ vanishes, rejection occurs with probability zero in each model; countability of words suffices.

Conversely choose a word with $a(w)>0$ and set
$e=\min\{1,a(w)/(2\kappa)\}>0$.
Accept every first herald except that, after precisely this word, reject with probability $e$ and then accept the next herald. The extra conditional cost after rejection is at most $\kappa$, from the actual post-herald state, by Lemma~\ref{lem:r67-first-herald}. Thus total cost in each model is at most
$t_\theta+e\kappa\lambda_\theta(w)\le N_0$.
A finite word matcher, a default accepting label and one post-rejection label implement this rule. Its only persistent information is those counted labels. This proves necessity and the realization assertion.
\end{proof}

\begin{corollary}[Active preparation coverage]\label{cor:r67-active-cover}
The rigidity condition holds if every $w\in\mathcal W$ has positive mass in at least one cost-active model $\theta$ with $t_\theta=N_0$. In particular it holds if one cost-active model has positive mass at every possible word. It also holds when the first-herald costs are equal across models. A positive definite herald effect and failure maps that send every nonzero positive input to a nonzero output give positive mass to every word, and therefore imply rigidity under \eqref{eq:r67-transience}.
\end{corollary}
\begin{proof}
An active positive-mass model makes the ratio in \eqref{eq:r67-margin} zero. For the last assertion, a finite failure composition has nonzero positive output, and pairing it with a positive definite herald effect gives positive trace. Lemma~\ref{lem:r67-first-herald} supplies an active model.
\end{proof}

The infimum in \eqref{eq:r67-margin} matters. Rigidity may also follow from models approaching the active set with both unused budget and word mass tending to zero, even if no active model gives that word positive mass. Conversely a word supported only where there is enough relative budget slack can permit boundary rejection. The theorem classifies stopping rigidity in this single-command uniformly transient class; it does not claim that a residual decision rule classifies multi-command boundary control.

\begin{example}[Active coverage is sufficient but not necessary]\label{ex:r67-flat-margin}
In dimension one take $\theta\in[0,1/2]$, $h(0)=0$, $h(\theta)=e^{-1/\theta}$ for $\theta>0$, and herald probability $p_\theta=1/(4-h(\theta))$. Give failures $1,0$ probabilities $(1-p_\theta)\theta$ and $(1-p_\theta)(1-\theta)$, with unit calls. Then $N_0=4$ and the only cost-active model is $0$, which gives no mass to words containing $1$. Nevertheless for a word of length $n$ containing $j$ occurrences of $1$,
\[
 \frac{N_0-t_\theta}{\lambda_\theta(w)}=
 \frac{e^{-1/\theta}}{p_\theta(1-p_\theta)^n\theta^j(1-\theta)^{n-j}}
 \longrightarrow0.
\]
Indeed the factors other than $\theta^j$ have positive limits, and $e^{-1/\theta}/\theta^j\le(j+1)!\theta$ by the exponential series. Words with no $1$ already have positive active-model mass. The exact criterion therefore gives boundary rigidity even though active support coverage fails. All primitive data are continuous and the herald probability is uniformly positive. No algebraic-data claim is made for this example.
\end{example}

\begin{example}[A paid rejection remains feasible at the robust least budget]\label{ex:r67-nonrigid}
Take two scalar models with unit costs. Model A has herald probability $1/4$ and only failure $0$; model B has herald probability $1/2$ and only failure $1$. The robust least budget is $4$. For the word consisting of one $1$, its mass is $1/4$ in B and zero in A, so $a(1)=8$. Accept the first herald except that after precisely this word reject with probability $1/2$, then accept the next herald. Its cost is $4$ in A and $2+(1/4)(1/2)2=9/4$ in B. A four-label matcher implements it. Thus a robust minimum-cost boundary need not eliminate stopping freedom merely because first-herald operation is modelwise cheapest.
\end{example}

\subsection{The entire robust optimal face, including ties}
Assume now that $a(w)=0$ on $\mathcal W$. At $N_0$ every feasible policy is a kernel $b:Y^*\to\Delta(C)$ with risk
\[
 \mathcal R_\theta(b)=\sum_{w,r}b(w,r)g_\theta(w,r).
\]
Its values on globally null words can be completed arbitrarily. Define $V=\min_b\sup_\theta\mathcal R_\theta(b)$. For a probability measure $\mu$ on $\Theta$, write $g_\mu(w,r)=\int g_\theta(w,r)\,d\mu(\theta)$ and
$A_\mu(w)=\operatorname*{argmin}_{r\in C}g_\mu(w,r)$.

\begin{theorem}[Exact robust optimal face]\label{thm:r67-optimal-face}
Primal and dual optima are attained, and
\begin{equation}\label{eq:r67-minimax}
 V=\max_{\mu\in\mathcal P(\Theta)}
                  \sum_w\min_{r\in C}g_\mu(w,r).
\end{equation}
For any attaining $\mu_*$, the complete set of robust optimal kernels is
\begin{equation}\label{eq:r67-face}
 \mathcal B_*=
 \left\{b:\ b(w,\cdot)\in\Delta(A_{\mu_*}(w))\ (w\in Y^*),\quad
                \mathcal R_\theta(b)\le V\ (\theta\in\Theta)\right\}.
\end{equation}
The global inequalities in \eqref{eq:r67-face} cannot in general be omitted. Wordwise Bayes ties do not authorize an arbitrary robust decoder.
\end{theorem}
\begin{proof}
The product of finite decision simplexes is compact convex. Every risk tail beyond length $T$ is at most $L\beta_T$, uniformly in the kernel and model. Finite sums are jointly continuous, so the full payoff is continuous and bilinear in the kernel and model measure. Compact minimax, or finite-prefix separation followed by compact limits with this uniform tail, gives \eqref{eq:r67-minimax} and both attainments \cite{Sion58}. The fixed-measure minimization separates wordwise. For an optimal kernel,
\[
 0\le\sum_{w,r}b(w,r)\bigl[g_{\mu_*}(w,r)-\min_u g_{\mu_*}(w,u)\bigr]
     =\int\mathcal R_\theta(b)\,d\mu_*-V\le0.
\]
Every summand is nonnegative, so all positive decoder weights are on minimizing actions. No strictly positive common word law is needed: zero-mass words simply impose no nontrivial local restriction. The all-model inequalities are necessary; together they imply optimality by the definition of $V$. This proves the characterization. The physical constructions in Theorem~\ref{thm:r67-embedding} below show why local ties alone are insufficient.
\end{proof}

\subsection{What replaces a finite deterministic residual quotient}
Let $\mathcal F=(\Delta(C))^{Y^*}$. Elements are full decoder-response functions, not hidden physical states. Define the affine shift $S_yf(z)=f(yz)$. A finite-dimensional polytope $P\subset\mathcal F$ is called invariant here when $S_yP\subset P$ for every failure letter. It is the convex hull of finitely many response functions; its extreme points are consequently finite. For $b\in P$ put
\[
 h(P,b)=|\operatorname{ext}P|+
                       \one_{\{b\notin\operatorname{ext}P\}}.
\]
The extra term records a prescribed pure initial label rather than an uncharged random initial state.

\begin{theorem}[Exact finite stochastic-memory criterion with robust ties]\label{thm:r67-polytope}
The least live-label count attaining $V$ at budget $N_0$ is exactly
\begin{equation}\label{eq:r67-polytope-count}
 \min\{h(P,b): b\in\mathcal B_*,\ b\in P\subset\mathcal F,
                    P\text{ a finite invariant polytope}\},
\end{equation}
with value infinity when no such pair exists. This equivalence includes stochastic transitions, randomized terminal decisions, impossible words, and all robust ties. It requires one common realization, not a different machine in each model.

For a deterministic kernel forced uniquely on every word, the criterion reduces to the number of distinct right residuals $z\mapsto b(wz)$, recovering the classical minimal output-machine count. Without uniqueness, a finite number of deterministic residual classes is neither the right requirement on an optimal randomized kernel nor a substitute for \eqref{eq:r67-polytope-count}.

If a realizing observer has $M$ live labels and $m$ distinct actual decoder values, event stopping uses at most $\max\{M,m\}$ total labels and state stopping at most $M+m$. These are the established packing counts, not extra live memory or a hidden clock.
\end{theorem}
\begin{proof}
An $M$-label first-herald observer has row-stochastic matrices $Q_y$ and decoder rows $D_i\in\Delta(C)$, from one prescribed initial label. Let $f_i(z)$ be its decision distribution on a future herald after failure word $z$, when initialized in label $i$. Then
\[
 f_i(\varnothing)=D_i,\qquad S_yf_i=\sum_j Q_y(i,j)f_j.
\]
Thus $P=\operatorname{conv}\{f_1,\ldots,f_M\}$ is invariant and contains its initial response $b$ as one listed point. It has at most $M$ extreme points. If $b$ is not extreme, it is itself a listed nonextreme point, so $h(P,b)\le M$ still holds. Its optimality puts $b$ in $\mathcal B_*$. This proves necessity and the lower count.

Conversely list the vertices of an admissible polytope, adding $b$ as an initial label only when it is not already a vertex. Every shift of each listed function belongs to $P$, and hence has fixed convex coefficients in those vertices. Choose these coefficients as its row of $Q_y$. Decode at a herald using that function's value at the empty word. Induction on word length proves that the machine's complete response from each label is the listed function. In particular its initial response is $b$. This produces exactly $h(P,b)$ labels, commits at the first herald and attains $V$. No physical state or posterior is read. Pruning unreachable labels cannot invalidate either implication or a minimal count.

When $b$ is deterministic, each distinct residual is an extreme point of the entire product $\mathcal F$: every coordinate is a vertex of $\Delta(C)$. Invariance puts every residual in $P$, so each must be one of its vertices. Conversely a finite set of distinct residuals is closed under the shifts and realizes the usual deterministic output machine. Its convex hull has precisely these vertices. This proves the reduction. The terminal counts follow by independently packing continuing and final roles as in Lemma~\ref{lem:roles}.
\end{proof}

Theorem~\ref{thm:r67-polytope} is an exact geometric existence criterion in response-function space. It is not a finite numerical test for arbitrary data, and it does not assume that a minimal stochastic realization lies in the affine span of the reached residuals. Such an assumption could discard valid hidden-state realizations. Positive realization by invariant cones or convex sets, and the distinction between linear rank and positive realization, have a long history \cite{Heller65,CarlylePaz71}. The new mathematical use here is the explicit least-budget condition, the robust tied optimal face, the pure-initial-label accounting, and the physical realization and obstruction results which follow.

\begin{proposition}[Fixed-memory attainment and boundary approximation]\label{prop:r67-density}
For every prescribed $M$, the least boundary risk among stochastic observers with at most $M$ live labels is attained. If \eqref{eq:r67-polytope-count} exceeds $M$, this risk is strictly larger than $V$. Nevertheless, a finite prefix of an optimal kernel, followed by one arbitrary decoder label, costs exactly $N_0$ in the robust sense and has excess risk at most $L\beta_T$. A sufficient live count is $1+\sum_{n=0}^{T-1}|Y|^n$. Modelwise cost equals $t_\theta$, not necessarily $N_0$ in every model.
\end{proposition}
\begin{proof}
Transition matrices and decoder rows at fixed $M$ range over compact simplexes. Finite-prefix risks are continuous and their uniform tail is at most $L\beta_T$. Thus worst-model risk is continuous and takes its minimum. Equality with $V$ is equivalent to an optimal realization by Theorem~\ref{thm:r67-polytope}. A prefix tree storing words shorter than $T$, followed by a common tail label, changes only terminal mass bounded by $\beta_T$ in every model. It still accepts the first herald and hence has the same modelwise costs $t_\theta$. This proves all assertions without taking an interior-budget completion theorem to zero slack.
\end{proof}

\section{Squared-prediction realizations of robust stochastic memory}\label{sec:robust-randomization}
A local Bayes tie can conceal global constraints on the entire randomized decision law. We construct such constraints using a fixed finite-dimensional instrument and a genuine squared terminal audit, not a time-indexed loss oracle. The construction gives exact finite-memory thresholds of every integer order, finite physical witnesses for strict smaller-memory gaps, and a finite-rank quantum law which has no finite stochastic realization. It also permits an exact algebraic fixed-memory optimization on a declared compact-model class.

\subsection{A physical embedding of a target channel response}
Fix a channel $\Psi$ on a $d_0$-dimensional protected system, a preparation $\rho$, and an effect $0\preccurlyeq E\preccurlyeq I$. Its target response is
\[
 u_n=\tr(E\Psi^n(\rho))\in[0,1],\qquad n\ge0.
\]
Choose $0<p<1$, $q=1-p$, $0<\bar t<1$, and $\varepsilon=1/4$. The unknown fixed model is $(t,\sigma)\in[0,\bar t]\times\{-1,1\}$. Add one protected classical two-state register, initially alive. Its channel $D_t$ discards coherences and maps the alive population to alive with probability $t$, otherwise to dead, while dead remains dead. Thus the real matrix entries of $D_t$ depend polynomially on $t$. A failure applies $D_t\otimes\Psi$ with probability $q$ and a herald applies the identity with probability $p$. There is one unit-cost command, one failure letter, and the fixed dictionary $C=\{-1,1\}$.

Let $A$ be the alive projection and set $T=2E-I$. Define two commuting Hermitian audit moments on the combined protected space by
\[
 v_\sigma=I-\sigma\varepsilon A\otimes T,
 \qquad m_\sigma=-\sigma\varepsilon(A\otimes I)/2.
\]
For $h=3/2$, use audit outcomes $z\in\{-h,0,h\}$ with effects
\begin{equation}\label{eq:r67-audit-povm}
 B_{\sigma,+}=\frac{v_\sigma+hm_\sigma}{2h^2},\qquad
 B_{\sigma,-}=\frac{v_\sigma-hm_\sigma}{2h^2},\qquad
 B_{\sigma,0}=I-v_\sigma/h^2.
\end{equation}
The observer commits before this audit is performed and pays loss $(z-r)^2$. Its conditional loss observable is
\begin{equation}\label{eq:r67-embedded-loss}
 F_{\sigma,r}=2I+\sigma\varepsilon A\otimes(rI-T).
\end{equation}

\begin{theorem}[A robust tied squared-task embedding]\label{thm:r67-embedding}
The displayed data define a legal prepared experiment with protected dimension $2d_0$, bounded squared loss, and least cost $1/p$. At that budget its robust risk is exactly $2$. Its unique optimal decision kernel is
\begin{equation}\label{eq:r67-target-kernel}
 b(n,+1)=u_n,\qquad b(n,-1)=1-u_n.
\end{equation}
For every fixed $t$, the equally weighted pair $(t,+1),(t,-1)$ is a least-favorable measure under which both actions are tied at every word. Nevertheless the all-model constraints force \eqref{eq:r67-target-kernel}, including its fractional values.

For an arbitrary first-herald kernel $b_n=b(n,+1)$, its exact worst-model excess is
\begin{equation}\label{eq:r67-generating-risk}
 \sup_{t,\sigma}\mathcal R_{t,\sigma}(b)-2
  =2\varepsilon p\sup_{0\le t\le\bar t}
                \left|\sum_{n\ge0}(qt)^n(b_n-u_n)\right|.
\end{equation}
Thus its exact finite stochastic-memory problem is the positive realization problem for the physically generated response $u_n$. Algebraic input channel, state, effect and constants give algebraic-polynomial primitive data.
\end{theorem}
\begin{proof}
One has $3I/4\preccurlyeq v_\sigma\preccurlyeq5I/4$ and $\|hm_\sigma\|\le3/16$. Consequently all three effects in \eqref{eq:r67-audit-povm} are positive and sum to $I$. Their first two moments are exactly $m_\sigma,v_\sigma$. Since $r^2=1$, the squared-loss observable is $v_\sigma-2rm_\sigma+I$, which is \eqref{eq:r67-embedded-loss}. The realized loss is bounded by $(h+1)^2=25/4$; its expectation is not an arbitrary signed score.

After $n$ failures the alive mass is $t^n$, independently of the target state $\Psi^n\rho$. The expected squared loss at decision $r$ is therefore
$2+\sigma\varepsilon t^n[r-(2u_n-1)]$.
First-herald commitment is forced at $1/p$, either by the indicator identity of Lemma~\ref{lem:boundary-herald} or by Theorem~\ref{thm:r67-rigidity}. Averaging the decision and summing its terminal mass $pq^n$ gives \eqref{eq:r67-generating-risk}. The kernel \eqref{eq:r67-target-kernel} makes every model's risk exactly $2$, while each symmetric pair of models averages to $2$ for every kernel. This proves the value and the stated least-favorable measures.

If a robust optimum differs at some coefficient, let $k$ be its first index with $b_k-u_k\ne0$. The power series in \eqref{eq:r67-generating-risk} vanishes on the interval by optimality. Divide it by $(qt)^k$ and let $t\downarrow0$. The remaining tail tends to zero because its coefficients are bounded, forcing $b_k-u_k=0$, a contradiction. This proves uniqueness without assuming an analytic inverse oracle. The target physical channel is not accessible to the observer except through the final audit. Its response must be reproduced by the observer's actual retained labels.
\end{proof}

This theorem explains why the local simplexes in \eqref{eq:r67-face} must be coupled by global risk inequalities. In this family every local simplex is the whole $\Delta(C)$ under the stated least-favorable measure, but there is only one robust optimal kernel.

The classical Erlang-chain rank mechanism is inherited from Theorem~\ref{thm:calibration-k}. The additions here are its realization by the same physical squared loss in every model and the finite physical witness theorem, which no longer needs continuum model uncertainty for the exact memory lower.

\subsection{Every finite exact memory order occurs}
For an integer $k\ge2$ and a rational $0<a<1$, take a classical target chain with $k-1$ transient stages and one absorbing dead state. At each transient stage it stays with probability $a$ and advances to the next stage with probability $1-a$; advancing from the last transient stage reaches dead. It starts at the first stage. Let $E$ indicate the transient states. Its survival response is
\begin{equation}\label{eq:r67-erlang-response}
 u_n=\sum_{j=0}^{k-2}\binom nj(1-a)^j a^{n-j},
\end{equation}
where a term is zero for $j>n$.

\begin{theorem}[A squared-loss realization of the exact memory hierarchy]\label{thm:r67-hierarchy}
The squared-task family of Theorem~\ref{thm:r67-embedding} built from this target has robust boundary value $2$ and least exact live-label count $k$. It has protected physical dimension $2k$. A $k$-label stochastic observer realizes the optimum at exactly $1/p$ cost. No observer with fewer live labels, even with arbitrary real stochastic transitions and randomized readouts, attains it. No deterministic full-history decision rule attains it either.

The full vector-valued Hankel response has rank exactly $k$, supplying the matching lower. Moreover, a finite collection of rational physical model parameters gives an explicit positive risk gap for every observer with fewer than $k$ live labels, as specified in Proposition~\ref{prop:r67-finite-witness}.
\end{theorem}
\begin{proof}
Use the target chain itself as the observer's failure-transition matrix; at a herald, decode $+1$ in transient stages and $-1$ in dead. Its probability of $+1$ after $n$ failures is \eqref{eq:r67-erlang-response}, so it attains the value with $k$ labels from a prescribed first label.

For any $M$-label stochastic observer with initial row $e_1$, failure matrix $Q$ and probability-of-$+1$ column $d$, its response is $b_n=e_1Q^nd$. The matrix with rows indexed by $i$ and columns by suffixes $j$ and decoder outcomes has rank at most $M$, since each entry factors through $e_1Q^i$. Include in particular the constant column $1$ and the columns $u_{i+j}$ for $0\le j\le k-2$, on rows $0\le i\le k-1$.

These $k$ columns are independent. Indeed $u_n=a^nP(n)$ for a polynomial $P$ of degree $k-2$ with nonzero leading coefficient. Its minimal shift annihilator is $(X-a)^{k-1}$: applying shift minus $a$ lowers the polynomial degree by exactly one, until zero. Equivalently, the finite confluent Vandermonde matrix with columns $1,a^i,ia^i,\ldots,i^{k-2}a^i$, $0\le i<k$, is nonsingular because $1\ne a$. A direct determinant proof is given in the supplement. The columns $u_{i+j}$ span the last $k-1$ columns because their translated degree-$k-2$ polynomials are independent. Hence the displayed Hankel block is nonsingular and every exact observer has at least $k$ labels. The construction supplies the upper.

At $n=k-1$, the survival probability lies strictly between zero and one: both no advances and all advances have positive probability. The unique optimal kernel is therefore fractional. A deterministic decision rule, even one with unlimited history, cannot equal it. The remaining assertion follows from the next proposition.
\end{proof}

\begin{proposition}[Finite physical witnesses for a smaller-memory gap]\label{prop:r67-finite-witness}
Let $u_n\in[0,1]$ be any target response for which the $k\times k$ matrix
\[
 A_u(i,0)=1,\qquad A_u(i,j+1)=u_{i+j},
 \quad 0\le i<k,\quad0\le j<k-1
\]
has nonzero determinant, where $k\ge2$. Put $K=2k-3$ and
$\Delta=|\det A_u|/[k!(k-1)]>0$.
Let $W_{jn}=(j/K)^n$ for $0\le j,n\le K$, taking $0^0=1$, and $C_0=\|W^{-1}\|_\infty$. Choose $0<r\le\min\{q\bar t,1/2,\Delta/(4C_0)\}$, put $Z_{jn}=(rj/K)^n$, and $C_r=\|Z^{-1}\|_\infty$.
Then every observer with fewer than $k$ live labels has worst-model risk at least
\begin{equation}\label{eq:r67-finite-gap}
 2+\varepsilon p\Delta/C_r.
\end{equation}
It suffices to take the finite models $t_j=rj/(Kq)$ and both signs $\sigma$. For rational target data and rational $p,\bar t$, all these constants and witnesses can be chosen rational. This is an explicit positive certificate, not an asserted sharp asymptotic memory rate.
\end{proposition}
\begin{proof}
The corresponding matrix $A_b$ for an $M<k$ label observer has rank at most $M$ and determinant zero. All entries lie in $[0,1]$. The Leibniz formula and telescoping a product of $k-1$ nonconstant entries give
$|\det A_u-\det A_b|\le k!(k-1)\max_{0\le n\le K}|u_n-b_n|$.
Thus one coefficient difference $e_n=b_n-u_n$ has magnitude at least $\Delta$.

Let $h(z)=\sum_{n\ge0}e_nz^n$. At $z_j=rj/K$, its tail after degree $K$ is at most $r^{K+1}/(1-r)$. Inverting the finite Vandermonde system gives
\[
 \max_{n\le K}|e_n|\le
 C_r\left(\max_j|h(z_j)|+\frac{r^{K+1}}{1-r}\right).
\]
Since $Z=W\operatorname{diag}(1,r,\ldots,r^K)$, one has $C_r\le C_0r^{-K}$ for $r\le1$. Our choice implies
$C_r r^{K+1}/(1-r)\le2C_0r\le\Delta/2$.
Consequently $\max_j|h(z_j)|\ge\Delta/(2C_r)$. Substitute this in the exact risk identity \eqref{eq:r67-generating-risk}, with both signs present at each node. The proof keeps the full infinite tail, not merely a fitted finite sample.
\end{proof}

\begin{corollary}[Finite-model rational exact-memory hierarchy]\label{cor:r67-finite-hierarchy}
For every integer $k\ge2$, there is a finite family of algebraic, in fact rational-superoperator, prepared classical instruments with the same fixed two-value dictionary and a genuine squared terminal audit, whose least-budget optimum is $2$ and whose exact least live-label count is $k$. A finite list of physical models gives an explicit strictly positive risk gap for every smaller stochastic observer. The first compulsory call and all subsequent calls remain charged. This finite-family statement does not assert uniqueness of the full-history optimal kernel.
\end{corollary}
\begin{proof}
Take rational $p,\bar t,a$ in the preceding hierarchy, and retain just the finite parameter list of Proposition~\ref{prop:r67-finite-witness}, with both signs at every node. Symmetry lower-bounds every policy's worst risk by $2$. The $k$-label target chain attains $2$ at every retained model. The finite-witness lower excludes every smaller memory. All maps and the finite audit effects have rational entries in the classical diagonal representation. First-herald commitment still has the common cost $1/p$, and the same predictable identity forces it at the boundary. The finite subfamily need not identify every coefficient of an arbitrary infinite decision kernel, which explains the final qualification.
\end{proof}

\subsection{The inherited rank obstruction under a common squared audit}
Theorem~\ref{thm:qubit-positive-phase} already establishes the rank-three obstruction for the signed calibration task. Here its same stochastic-realization obstruction is transported, with full proof retained, to one common squared-loss audit. The obstruction itself is not a newly discovered fact.

\begin{theorem}[A rank-three protected quantum response with no finite realization]\label{thm:r67-rank-three}
Take the target qubit with
\[
 \rho=E=|0\rangle\langle0|,\qquad
 U=\frac1{\sqrt5}\begin{pmatrix}2&-1\\1&2\end{pmatrix},
 \qquad\Psi(A)=UAU^*.
\]
Its response is $u_n=(1+\cos(n\vartheta))/2$, where $e^{i\vartheta}=(3+4i)/5$. The full two-outcome Hankel response has rank three, but no finite stochastic observer realizes it exactly.

The embedding theorem therefore gives a compact uncertain-model squared-prediction family on protected dimension four whose unique full-history boundary optimum is randomized, whose least-favorable wordwise actions are all tied, and whose exact live-memory count is infinite. At each finite memory bound its optimal risk is strictly above $2$. A prefix of length $T$, followed by one tail label, has excess risk at most
\begin{equation}\label{eq:r67-rank-three-tail}
 \frac{2\varepsilon p(q\bar t)^T}{1-q\bar t}
\end{equation}
and costs exactly $1/p$.
\end{theorem}
\begin{proof}
The response formula follows by the direct qubit rotation. Its three distinct exponential modes are $1,\zeta^n,\bar\zeta^n$, where $\zeta=(3+4i)/5$. Their nonzero coefficients and the Vandermonde determinant give Hankel rank three. The number $\zeta$ is not a root of unity: otherwise $\zeta+\bar\zeta=6/5$ would be a rational algebraic integer and hence an integer.

If a finite stochastic matrix $Q$ realized $u_n=e_1Q^nd$, its generating function would be $e_1(I-zQ)^{-1}d$. But
\[
 \sum_{n\ge0}u_nz^n=
 \frac1{2(1-z)}+\frac1{4(1-\zeta z)}+
                         \frac1{4(1-\bar\zeta z)}.
\]
The pole at $z=\zeta^{-1}$ forces $\zeta$ to be an eigenvalue of $Q$. Every unit-modulus eigenvalue of a finite row-stochastic matrix is a root of unity. To verify the needed fact, choose a complex eigenvector and a coordinate of maximal nonzero modulus. Equality in the stochastic triangle inequality forces all positive successors to have the same modulus and phase multiplied by that eigenvalue. Following positive edges reaches a directed cycle; the phase must return after its finite length. This gives a root of unity, a contradiction.

The embedding gives uniqueness of the optimal randomized kernel, and Proposition~\ref{prop:r67-density} gives strict attained gaps at fixed finite memory. Matching $u_n$ for $n<T$ and assigning any tail probability changes each coefficient by at most one. Summing the tail in \eqref{eq:r67-generating-risk} gives \eqref{eq:r67-rank-three-tail}. There is no change to first-herald stopping or its paid cost.
\end{proof}

Irrational quantum rotations and the separation of linear from positive stochastic realizations are established mechanisms, not new priority claims \cite{Heller65,CarlylePaz71}. The assertion here concerns their exact robust squared-prediction embedding, all-tied least-favorable actions, fixed-dictionary boundary optimality and real-valued stochastic-controller exclusion. It does not solve the sharp intermediate-memory risk curve for the different noisy-qubit sign task in Theorem~\ref{thm:two-label-quantum}.

\subsection{Exact optimization at any prescribed finite memory}
The invariant-polytope criterion is not asserted effectively decidable in general. For a specified memory bound, however, the physical instrument gives a finite algebraic optimization without an unproved exact posterior algorithm.

\begin{theorem}[Algebraic fixed-memory boundary optimization]\label{thm:r67-algebraic-memory}
Suppose $\Theta$ is a nonempty compact semialgebraic set over the real algebraic numbers, and the real Hermitian-coordinate entries of all prepared states, instruments, costs and loss observables are algebraic-polynomial functions of its parameters. Assume \eqref{eq:r67-transience} and boundary rigidity. For each prescribed live bound $M$, the exact robust boundary optimum $V_M$ is an attained real algebraic number and is computable by real-algebraic elimination. An attaining stochastic controller can be chosen with real algebraic entries.

A completely primitive, decidable sufficient subclass for rigidity is
$\mathcal I_{\theta,s}^*I\succ0$ and $\mathcal I_{\theta,y}^*I\succ0$ for every $\theta,y$; compactness then also gives uniform first-step transience. These conditions allow model- and state-dependent report probabilities. Both squared-task embeddings above belong to this subclass.

For the same algebraic data, $V$ is a computable real via finite-prefix games with the explicit $L\beta_T$ tail. Whenever $V_M>V$, finite computations eventually certify that strict inequality. No general terminating test for the equality $V_M=V$, no test that some finite $M$ attains $V$, and no polynomial complexity for quantifier elimination is inferred.
\end{theorem}
\begin{proof}
Write the controller variables as row-stochastic $Q_y$ and decoder probabilities $D_i(r)$, with initial label 1. For each model its exact score potentials are the unique Hermitian solution of
\begin{equation}\label{eq:r67-score-resolvent}
 K_i=\mathcal I_{\theta,s}^*\!\left(\sum_rD_i(r)F_{\theta,r}\right)
       +\sum_y\mathcal I_{\theta,y}^*\!\left(\sum_jQ_y(i,j)K_j\right),
       \qquad 1\le i\le M.
\end{equation}
Indeed the associated positive continuation map on block observables sends the identity tuple, at its $n$th power, to the tuple with common entry $(\mathcal T_\theta^*)^nI$. Its norm on Hermitian tuples is therefore at most $\beta_n$. Its Neumann series converges, giving uniqueness and $0\preccurlyeq K_i\preccurlyeq LI$. The risk is exactly $\tr(\rho_\theta K_1)$. This proves the equation's meaning and boundedness from the raw instrument, not by assuming a certificate exists.

For a proposed upper $v$, impose stochastic row constraints and the statement
\[
 \forall\theta\in\Theta\ \exists(K_1,\ldots,K_M):\quad
 \eqref{eq:r67-score-resolvent},\quad
 0\preccurlyeq K_i\preccurlyeq LI,\quad
 \tr(\rho_\theta K_1)\le v.
\]
In real coordinates these are polynomial equations and inequalities; positive semidefiniteness is expressed by principal minors. Projection and quantifier elimination over the real algebraic field produce the exact feasible value set \cite{Seidenberg54}. Compact controller space and the uniform tail prove its lower endpoint is attained. A finite semialgebraic endpoint over that field is algebraic, and the nonempty optimal-controller fiber contains an algebraic point. This gives the value and a controller by a finite elimination procedure. The $Md^2$ modelwise potential coordinates are analytic variables in this computation, not runtime observer memory.

Positive definiteness of the primitive effects is a finite quantified algebraic condition. It implies positive mass at every word, and compactness bounds the herald effect below by a common positive multiple of $I$. Corollary~\ref{cor:r67-active-cover} therefore gives rigidity. The general theorem's rigidity premise is not claimed algorithmically checked merely by writing the infinite-word criterion.

Finally, truncate the decision kernel to words shorter than $T$. The finite-prefix robust minimum is an algebraic quantified problem of the same kind, now with one decision row per word. Its value is below $V$ and within $L\beta_T$ of it because any prefix can be completed by a first-herald tail decoder. A rational positive lower certificate for $\delta$ can be found by increasing-precision algebraic tests, so these are effective converging intervals. Once the exact algebraic value $V_M$ is strictly above $V$, some finite upper interval endpoint for $V$ is strictly below $V_M$. This certifies nonattainment at that fixed bound. In the equality case that search need not terminate.
\end{proof}

\section{The causal task-order transfer theorem}\label{sec:central}
Suppose that every terminating state-stopped competitor reaches at least $B_0\ge1$ live labels under $\theta_*$. Such a certificate must concern actual physical command compatibility, not printed action names. Define
\begin{equation}\label{eq:widths}
\begin{array}{c|cc}
 i&m_-(M)&m_+(M)\\\hline
 \cp&M&\lfloor M/D_0\rfloor\\
 \ev&M&\lfloor M/D_0\rfloor\\
 \st&M-B_0&\lfloor(M-B)/D_0\rfloor .
\end{array}
\end{equation}
The event upper also requires $M\ge B$. The trivial state-stopped lower $B_0=1$ is valid for a procedure initially waiting to make a call.

\begin{theorem}[Affine task-order transfer with a paid reference and common construction]\label{thm:transfer}
Assume all the prepared-kernel, full-report, task-order, effort, occupation, resource and common-construction hypotheses of Section~\ref{sec:data}. Fix $N,M$, put $b_-=N/\kappa_-$, $b_+=N/\kappa_+$, and suppose $b_-\ge1$ and $m_-(M)\ge1$. Set
\[
 t_N=(1-A_N)_+/u_-.
\]
If $t_N>0$, every one of the three risks satisfies
\begin{equation}\label{eq:central-lower}
 R_i^\Theta(N,M)\ge
 \bigl[l_-t_N\Q_{m_-(M)}(b_-/t_N)-E_N\bigr]_+.
\end{equation}
If $t_N=0$, the asserted lower is zero, with no evaluation of $\Q$ at infinity. An empty legal class still has risk $+\infty$.

Suppose also $m_+(M)\ge1$, $b_+\ge1$, and, for event stopping, $M\ge B$. If
\begin{equation}\label{eq:beta}
 \beta_N=\frac{l_+b_+}{1+b_+a_+}\ge1,
\end{equation}
one common Borel program with actual expected cost at most $N$ in every model attains
\begin{equation}\label{eq:central-upper}
 R_i^\Theta(N,M)\le
 \frac{u_+(1+b_+a_+)}{l_+}\Q_{m_+(M)}(\beta_N)+b_+e_+.
\end{equation}
It commits almost surely. If $\beta_N<1$, the always-accept first-attempt construction instead gives $R_i^\Theta(N,M)\le u_+\Q_{m_+(M)}(1)+e_+$ at cost at most $\kappa_+\le N$.

The lower is against every legal competing strategy and uses one actual witness. The upper is one parameter-independent strategy for all models. The two effort constants need not agree. Neither direction assumes that the comparison experiment belongs to $\Theta$ or that a normalized predictive update is uniformly contracting.
\end{theorem}

\begin{corollary}[Scale-uniform relative transfer]\label{cor:relative}
Under Theorem~\ref{thm:transfer}, suppose $A_N=E_N=a_+=e_+=0$. If $l_+b_+\ge1$, then
\begin{equation}\label{eq:relative}
 \frac{l_-}{u_-}\Q_{m_-(M)}(u_-b_-)
 \le R_i^\Theta(N,M)
 \le\frac{u_+}{l_+}\Q_{m_+(M)}(l_+b_+).
\end{equation}
The factors $l_\pm,u_\pm$ may be fixed nonunit constants, independent of budget and resolution. In particular no term proportional to $b_\pm$ times a fixed absolute calibration error is required in this relative-order class.
\end{corollary}
\begin{proof}
Substitute $t_N=1/u_-$ and $\beta_N=l_+b_+$ in the theorem.
\end{proof}

\begin{corollary}[Separate exact checkpoint and autonomous laws]\label{cor:exact}
Suppose the certificates are exact: $l_\pm=u_\pm=1$, all defects vanish, $\kappa_-=\kappa_+=\kappa$, and exports are the identity. Then, for $N\ge\kappa$,
\[
 R_\cp^\Theta(N,m)=\Q_m(N/\kappa).
\]
The constructed autonomous upper uses $B+m$ state-stopped labels, or $\max(B,m)$ event-stopped labels. If the mandatory live count is sharp, $B_0=B$, then
\begin{align}
 R_\st^\Theta(N,B+m)&=\Q_m(N/\kappa)\quad(m\ge1),\label{eq:typed-exact}\\
 R_\ev^\Theta(N,M)&=\Q_M(N/\kappa)\quad(M\ge B).\label{eq:event-exact}
\end{align}
The subtraction by $B$ belongs only to the state-stopped model. Nonidentity exports require the expanded dictionary in \eqref{eq:widths}, unless an optimal dictionary is invariant under all exports. These exact stopping distinctions are inherited from the complete companion and are reproved here to fix the resource semantics.
\end{corollary}
\begin{proof}
The same lower and upper profile coincide in each asserted interface. Identity exports give $D_0=1$. The packing and converse widths are proved in Lemma~\ref{lem:roles}. No lower on live states is used to forbid event-stopped reuse.
\end{proof}

The familiar absolute-defect theorem is also a specialization: take $l_\pm=u_\pm=1$ while retaining $A_N,E_N,a_+,e_+$. Thus relative and absolute calibration are composed in one statement. The chosen metric is the common scored task throughout; it is not changed between the converse and the construction.

\section{Proof of the central theorem}
\begin{lemma}[Profile attainment]\label{lem:attainment}
For every finite $m\ge1$ and $b\ge1$, the minimum in \eqref{eq:profile} is attained. The profile is nonincreasing in $m,b$ and left-continuous in $b>1$.
\end{lemma}
\begin{proof}
Write a dictionary as an ordered tuple in the compact set $H^m$, permitting repetitions. The feasible densities are a weak-star closed subset of the radius-$b$ ball of $L^\infty(\nu)$: positivity, the upper bound and the mass equality are each closed under pairing with $L^1(\nu)$. They are therefore weak-star compact. As the tuple varies, its squared-distance function varies uniformly on $H$, since all points belong to a bounded set. Along a convergent product net, split the objective difference into this uniform variation and the weak-star pairing against the fixed bounded measurable loss at the limit tuple. This proves continuity of the objective and attainment. Nets avoid an unwarranted sequential compactness assumption. Monotonicity follows from inclusion of feasible sets. For left continuity take an optimizer $h$ at $b>1$ and let $h_\lambda=(1-\lambda)h+\lambda$. It has unit mass and cap $b_\lambda=(1-\lambda)b+\lambda<b$. Its objective tends to the optimum as $\lambda\downarrow0$, while monotonicity bounds the optimum at $b_\lambda$ between the optimum at $b$ and this objective.
\end{proof}

\begin{lemma}[Actual-history comparison measure]\label{lem:comparison}
Fix any robustly admissible competitor of cost at most $N$, with its fixed terminal dictionary $C$. There are finite nonnegative measures $\xi_c$ on $X$, $c\in C$, such that $\xi=\sum_c\xi_c$ obeys
\begin{align}
 \xi&\le (\E_*A_\tau)\nu\le b_-\nu,\label{eq:measure-cap}\\
 1&\le u_-\xi(X)+A_N,\label{eq:measure-mass}\\
 \E_*\|Z-\widehat r\|^2
 &\ge l_-\sum_{c\in C}\int d_c\,d\xi_c-E_N.
 \label{eq:measure-risk}
\end{align}
The measure $\xi$ is not asserted to be the terminal law of any comparison execution.
\end{lemma}
\begin{proof}
Adjoin an independent uniform random variable to the current report on both sides. This is current-call randomness, not a seed retained between calls; it is sufficient to represent the competitor's randomized output and stopping decision. A finite-output kernel has such a representation by cumulative probabilities. At a fixed reached productive history $h$ and chosen command $a$, the post-report decision then gives an acceptance indicator $q_h(y)$ and an output $r_h(y)\in C$ on acceptance. Give $r_h$ an arbitrary fixed value on rejection. The notation $y$ now includes the current independent random variable. Integrating the original certificates over that variable preserves their constants. Dependence on the entire earlier history is permitted in these functions.

For each Borel $D\subset X$, set
\begin{equation}\label{eq:xi-definition}
 \xi_c(D)=\E_*\sum_{j\le\tau:\ j\ {
m productive}}
 \int\one_{\{\pi_{H_{j-1},a_j}(y)\in D\}}
 q_{H_{j-1}}(y)\one_{\{r_{H_{j-1}}(y)=c\}}
 P^0_{H_{j-1},a_j}(dy,H).
\end{equation}
The outer expectation is the actual pre-call occupation at $\theta_*$. All terms are jointly measurable and nonnegative. Tonelli defines the measures, whose total is finite because $\E_*A_\tau\le b_-$. Summing over $c$ and dropping the acceptance indicator gives \eqref{eq:measure-cap}: each comparison row has mark law $\nu$.

The analogous sum with the actual rows instead of $P^0$ has mass exactly one, because the actual procedure commits exactly once almost surely. Apply \eqref{eq:lower-mass} at each actual reached history, then sum. The occupation certificate yields \eqref{eq:measure-mass}. Apply \eqref{eq:lower-task} in the same way to the disjoint actual commitment events. Their expected losses sum to the actual risk. The comparison contribution is $d_{r_h(y)}(\pi_{h,a}y)$ by \eqref{eq:full-moments}; this identity conditions on the full current report and therefore does not overlook extra information seen by the competitor. The accumulated loss defect is at most $E_N$, giving \eqref{eq:measure-risk}. One may first truncate the number of calls. Monotone convergence for the nonnegative contributions and the finite defect bounds justify the limit.

No actual past is replaced by a comparison past. In particular the competitor need not terminate, or obey a cost bound, if hypothetically run under the comparison experiment. This distinction is why the argument applies when that experiment is not an operating model.
\end{proof}

\begin{remark}[Why the full report is required]
A constant mark can conceal a visible fair bit which also equals the future audit. A comparison using only that mark has $f=0,w=1$, but a two-value decoder reading the bit predicts it with zero loss. The hypotheses therefore condition on the full report seen by the competitor; matching a compressed mark marginal alone cannot establish the lower.
\end{remark}

\begin{lemma}[Live roles, output roles and physical compatibility]\label{lem:roles}
A program with $\ell$ live roles and $t$ output roles has a state-stopped implementation with $\ell+t$ labels and an event-stopped implementation with $\max(\ell,t)$ labels, preserving every path cost and complete stopped scored law. Conversely an event-stopped $M$-label program has at most $M$ decoder values. A state-stopped $M$-label program with at least $B_0$ reached live labels has at most $M-B_0$ terminal decoder values.

For a physical live-state certificate, consider finitely many mandatory pre-call contexts $I$ with no free context input. Suppose the commands legal at a history of context $i$ are exactly $\mathcal U_i$, and every terminating competitor visits every context with positive probability under the same witness. This positive probability may depend on the competitor; no uniform lower bound over competitors is required. Let $b_{\rm phys}$ be the least size of a cover of $I$ by subsets $J$ with $\bigcap_{i\in J}\mathcal U_i\ne\varnothing$. Then at least $b_{\rm phys}$ live labels are required. This certificate is invariant under duplicate names and bijective stateless recodings of the same physical commands.
\end{lemma}
\begin{proof}
For state stopping take disjoint copies of the live and terminal role sets and type them zero and one. For event stopping embed the live and terminal role sets separately into one alphabet of size $\max(\ell,t)$. On continuing transitions use the live embedding; on commitment emit the final event and use the terminal embedding. A command row is read only on a continuing branch, while the decoder is read only after commitment. Sharing their label therefore causes no conflict and uses no continuing mode bit. Couple the original and recoded coins. Every finite stopped prefix and every path cost agree, and almost-sure termination gives equality of the stopped scored laws. A fixed decoder has one value per retained label. In state stopping the global halt map excludes each reached live label from the output set, proving its subtraction. Event stopping has no such exclusion.

A live label reached in contexts $J$ must have one command distribution supported on all the corresponding $\mathcal U_i$. Since $J$ is finite, the union of the exceptional null command sets is null; hence their intersection supports a probability and is nonempty. The sets of contexts served by the labels cover $I$, proving the bound. Conversely a nonempty intersection contains a command and hence a Dirac command distribution. Replacing names without changing executed kernels, legal domains, costs and reports does not change these intersections. A translator which itself remembers a hidden phase is not such a recoding; its memory must be charged.
\end{proof}

\begin{proof}[Proof of Theorem~\ref{thm:transfer}]
Fix a competing program. Lemma~\ref{lem:roles} bounds its decoder width by $m_-(M)$. Let $s=\xi(X)$ in Lemma~\ref{lem:comparison}. If $t_N>0$, then \eqref{eq:measure-mass} gives $s\ge t_N$. Scale the measure to $\xi'=(t_N/s)\xi$. It has mass $t_N$ and remains bounded above by $b_-\nu$. Thus $h=(d\xi'/d\nu)/t_N$ has unit mass and cap $b_-/t_N$. Nearest-dictionary projection and nonnegativity imply
\[
 \sum_{c\in C}\int d_c\,d\xi_c
 \ge\int [w+\dist(f,C)^2]\,d\xi
 \ge t_N\Q_{m_-(M)}(b_-/t_N).
\]
Together with \eqref{eq:measure-risk} this proves the asserted lower for this competitor. Its worst-model risk dominates its risk at $\theta_*$. Taking the infimum over competitors gives \eqref{eq:central-lower}, without exchanging an infimum and a supremum. Nonnegative loss gives the zero truncation and the stated case $t_N=0$.

For the upper assume \eqref{eq:beta}. Choose an attained optimizer $(C,h)$ at width $m_+(M)$ and cap $\beta_N$. Let $r_C(x)$ be the first nearest point to $f(x)$ in a fixed ordering of $C$. After each attempt of $\mathcal G$, accept $x$ with probability $h(x)/\beta_N$, export $T_Jr_C(x)$ and commit; otherwise return to the start routing label. All these operations are Borel. The choice $J$ is used within the current attempt; any part retained thereafter is charged through the live role or exported decoder, not a free seed. The exports are assumed jointly measurable and take at most $D_0|C|$ distinct values, so their finite decoder and transitions are Borel as well.

Conditional on every surviving attempt history, the acceptance probability is at least
\begin{equation}\label{eq:hazard}
 p=\frac{l_+}{\beta_N}-a_+=\frac1{b_+}>0.
\end{equation}
Consequently the probability of reaching attempt $j+1$ is at most $(1-p)^j$, by induction on conditional survival probabilities. In particular commitment is almost sure and the expected number of attempts is at most $1/p=b_+$. Each begun attempt has conditional expected cost at most $\kappa_+$, not merely a cost conditional on a favorable outcome. Summing the predictable attempt costs gives actual expected cost at most $\kappa_+b_+=N$.

The conditional unnormalized loss contributed by an attempt is at most $u_+\Q_{m_+(M)}(\beta_N)/\beta_N+e_+$. Sum its contributions on the disjoint terminal events, using the same survival bound. The result is at most
\[
 b_+\left[\frac{u_+}{\beta_N}\Q_{m_+(M)}(\beta_N)+e_+\right],
\]
which is \eqref{eq:central-upper}. The program has $B$ live roles and at most $D_0m_+(M)$ terminal roles. Lemma~\ref{lem:roles} gives exactly the sufficient widths in \eqref{eq:widths}. A checkpoint counts only the terminal width. No physical renewal law was used: the conditional hazard controls survival even if the hidden state after rejection differs arbitrarily from its former state.

If $\beta_N<1$, commit after the first attempt, and use an optimizer at cap one. The loss certificate with $q=1$ gives $u_+\Q_{m_+(M)}(1)+e_+$, and the whole-attempt cost is at most $\kappa_+$. This proves all remaining assertions.
\end{proof}

\section{Scope and continuing mathematical questions}\label{sec:r70-scope}
The inherited boundary criterion concerns one globally legal command, finitely many reports, compact continuous model uncertainty, a fixed finite decoder dictionary, and primitive uniform transience. It permits non-scalar herald effects, noncommuting channels and model-dependent word laws. It does not classify boundary feedback with several commands. Where the normalized unused-budget condition fails, the theorem supplies an actual rejecting policy but does not replace its remaining control problem by a decision quotient.

Robust ties are classified by a finite positive stochastic realization of an optimal kernel, not by arbitrary selection from wordwise Bayes ties. The invariant-polytope criterion is complete within the stated rigid class, but is not a general algorithm. Exact risk at each prescribed algebraic memory bound is computable; deciding whether some finite memory attains the unrestricted optimum is a different question. The rational exact-memory hierarchy and rank-three counterexample are physically realized squared-prediction constructions, not claims to settle a separately recognized open quantum-control problem. The new exact instrument simulator keeps its latent state hidden and does not itself construct an optimal finite-memory observer. Its state-positive noise exponent does not bound arbitrary response realizations. The present exact-response theorem has a different proof: its worst-phase label law is exact, and its fixed badly-approximable phase law has matching exponents. The stronger logarithmic lower at the older algebraic angle is not asserted sharp. The general-$M$ squared-sign risk frontier and the intermediate phase-type timing frontier remain open here.

The geometric transfer still requires its own actual-history converse and common paid construction. Its full finite-width, rank, singular-range, regular, mixed and atomic consequences remain in the current supplement, together with nonintervening sensing ancestry and the fixed-law predictive quotient. The independent ordered-measurement, two-stage reset communication and rank-constrained discrimination sources remain preserved at their immutable identities; they are not inferred from the new protected-channel embedding. General unknown-model predictive realization and matching optimal program, precision, workspace and time bounds remain independent research targets.
The cubic-root arithmetic hypothesis is global at every denominator. It is proved here for a quadratic irrational phase, but not for the distinct rotation with $e^{i\vartheta}=(3+4i)/5$. No general matched scalar-memory rate for every irrational phase or every multi-frequency quantum response is asserted. Algebraic fixed-size spectral certificates and ideal exact label counts also do not supply matched program, precision, sampling or working-storage bounds.
The leakage theorem concerns vanishing probability outputs and a bounded stochastic decoder with a prescribed initial label. Its exact extra label is a continuing zero-output state, not an unobserved commitment event. The finite-model witnesses grow in number with the memory bound under test; they are not one fixed finite family excluding all finite controllers. Positive background gives a finite exact construction but not a sharp background-dependent lower rate. The strict trace-contraction example approaches a boundary stationary state and therefore does not contradict the state-interior simulator theorem. The new results do not replace the unresolved general multi-command boundary problem or the older squared-sign and timing risk frontiers.
\begin{thebibliography}{99}
\bibitem{Blackwell} D. Blackwell, \emph{Equivalent comparisons of experiments}, Ann. Math. Statist. \textbf{24} (1953), 265--272. DOI: 10.1214/aoms/1177729032.
\bibitem{CGM97} J. A. Cuesta-Albertos, A. Gordaliza and C. Matr\'an, \emph{Trimmed $k$-means: an attempt to robustify quantizers}, Ann. Statist. \textbf{25} (1997), 553--576. DOI: 10.1214/aos/1031833664.
\bibitem{GG99} L. A. Garc\'ia-Escudero and A. Gordaliza, \emph{Robustness properties of $k$ means and trimmed $k$ means}, J. Amer. Statist. Assoc. \textbf{94} (1999), 956--969. DOI: 10.1080/01621459.1999.10474200.
\bibitem{SC} C. R. Shalizi and J. P. Crutchfield, \emph{Computational mechanics: pattern and prediction, structure and simplicity}, J. Statist. Phys. \textbf{104} (2001), 817--879. DOI: 10.1023/A:1010388907793.
\bibitem{BC} N. Barnett and J. P. Crutchfield, \emph{Computational mechanics of input--output processes: structured transformations and the $\varepsilon$-transducer}, J. Statist. Phys. \textbf{161} (2015), 404--451; arXiv:1412.2690.
\bibitem{AMTW} A. Ambainis, M. Mosca, A. Tapp and R. de Wolf, \emph{Private quantum channels}, Proc. 41st Annual Symposium on Foundations of Computer Science (2000), 547--553. DOI: 10.1109/SFCS.2000.892142.
\bibitem{Buscemi} F. Buscemi, \emph{Comparison of quantum statistical models: equivalent conditions for sufficiency}, Comm. Math. Phys. \textbf{310} (2012), 625--647. DOI: 10.1007/s00220-012-1421-3.
\bibitem{BR14} N. Bertschinger and J. Rauh, \emph{The Blackwell relation defines no lattice}, Proc. IEEE International Symposium on Information Theory (2014), 2479--2483. DOI: 10.1109/ISIT.2014.6875280; arXiv:1401.3146.
\bibitem{Banerjee05} A. Banerjee, X. Guo and H. Wang, \emph{On the optimality of conditional expectation as a Bregman predictor}, IEEE Trans. Inform. Theory \textbf{51} (2005), 2664--2669. DOI: 10.1109/TIT.2005.850145.
\bibitem{Stengel96} B. von Stengel, \emph{Efficient computation of behavior strategies}, Games Econom. Behav. \textbf{14} (1996), 220--246. DOI: 10.1006/game.1996.0050.
\bibitem{Balder89} E. J. Balder, \emph{On compactness of the space of policies in stochastic dynamic programming}, Stochastic Process. Appl. \textbf{32} (1989), 141--150. DOI: 10.1016/0304-4149(89)90058-6.
\bibitem{Barry14} J. Barry, D. T. Barry and S. Aaronson, \emph{Quantum partially observable Markov decision processes}, Phys. Rev. A \textbf{90} (2014), 032311. DOI: 10.1103/PhysRevA.90.032311; arXiv:1406.2858.
\bibitem{Ying18} S. Ying and M. Ying, \emph{Reachability analysis of quantum Markov decision processes}, Inform. Comput. \textbf{263} (2018), 31--51. DOI: 10.1016/j.ic.2018.09.001; arXiv:1406.6146.
\bibitem{Saldi25} N. Saldi, S. Sanjari and S. Y\"uksel, \emph{Quantum Markov decision processes: general theory, approximations, and classes of policies}, SIAM J. Control Optim. \textbf{63} (2025), 1156--1183. DOI: 10.1137/24M1643864; arXiv:2402.14649.
\bibitem{Seidenberg54} A. Seidenberg, \emph{A new decision method for elementary algebra}, Ann. of Math. (2) \textbf{60} (1954), 365--374. DOI: 10.2307/1969640.
\bibitem{Chatterjee15} K. Chatterjee, M. Chmel\'ik, R. Gupta and A. Kanodia, \emph{Optimal cost almost-sure reachability in POMDPs}, Proc. AAAI Conference on Artificial Intelligence \textbf{29} (2015). DOI: 10.1609/aaai.v29i1.9683.
\bibitem{Andriushchenko23} R. Andriushchenko, A. Bork, M. \v{C}e\v{s}ka, S. Junges, J.-P. Katoen and M. Mac\'ak, \emph{Search and explore: symbiotic policy synthesis in POMDPs}, Computer Aided Verification (2023); arXiv:2305.14149v2.
\bibitem{Halmos69} P. R. Halmos, \emph{Two subspaces}, Trans. Amer. Math. Soc. \textbf{144} (1969), 381--389. DOI: 10.1090/S0002-9947-1969-0251519-5.
\bibitem{Kessler21} D. A. Kessler, E. Barkai and K. Ziegler, \emph{First-detection time of a quantum state under random probing}, Phys. Rev. A \textbf{103} (2021), 022222. DOI: 10.1103/PhysRevA.103.022222; arXiv:2012.01763.
\bibitem{DelVecchio26} G. Del Vecchio Del Vecchio and S. N. Majumdar, \emph{Optimal detection of quantum states via projective measurements}, J. Phys. A: Math. Theor. \textbf{59} (2026), 035001. DOI: 10.1088/1751-8121/ae34fe; arXiv:2509.08556.
\bibitem{Companion} Q. Qi, \emph{General Theta Foundations I: Acquired Geometry and Causal Resource Transfer}, complete technical companion on absolute calibration, exact stopping resources and nonintervening sensing, 2026. The complete retained R58 native source and its independently built companion accompany this manuscript; their identities are fixed in the source-provenance and preservation records.
\bibitem{Tzeng92} W.-G. Tzeng, \emph{A polynomial-time algorithm for the equivalence of probabilistic automata}, SIAM J. Comput. \textbf{21} (1992), 216--227. DOI: 10.1137/0221017.
\bibitem{LiYuYing12} Y. Li, N. Yu and M. Ying, \emph{Termination of nondeterministic quantum programs}, arXiv:1201.0891 (2012).
\bibitem{Fu24} J. Fu, H. Jiang, M. Xu, Y. Deng and Z.-B. Li, \emph{Algorithmic analysis of termination problems for nondeterministic quantum programs}, arXiv:2402.15827 (2024).

\bibitem{LiuVerdu18} J. Liu and S. Verd\'u, \emph{Rejection sampling and noncausal sampling under moment constraints}, Proc. IEEE International Symposium on Information Theory (2018), 1565--1569. DOI: 10.1109/ISIT.2018.8437857.
\bibitem{GocFlamich24} D. Goc and G. Flamich, \emph{On channel simulation with causal rejection samplers}, Proc. IEEE International Symposium on Information Theory (2024). DOI: 10.1109/ISIT57864.2024.10619339; examined text arXiv:2401.16579v2 (3 May 2024).
\bibitem{LiLiu26} Z. Li and T. Liu, \emph{Ranking experiments under sequential sampling}, arXiv:2608.19897v2 (28 August 2026). Version 1 appeared under the title \emph{The order of binary experiments under endogenous stopping}, with Z. Li as sole author; the version-2 statement is the one compared here.
\bibitem{Ying14pre} S. Ying and M. Ying, \emph{Reachability analysis of quantum Markov decision processes}, arXiv:1406.6146 (2014), Theorem 3.6; published journal version listed above as \cite{Ying18}.

\bibitem{KaraYuksel20} A. D. Kara and S. Y\"uksel, \emph{Robustness to incorrect system models in stochastic control}, SIAM J. Control Optim. \textbf{58} (2020), 1144--1182. DOI: 10.1137/18M1208058; examined text arXiv:1803.06046v7.
\bibitem{Demirci25} Y. E. Demirci, A. D. Kara and S. Y\"uksel, \emph{Sensitivity of filter kernels and robustness bounds to transition and measurement kernel perturbations in partially observable stochastic control}, arXiv:2508.10658v3 (2025).
\bibitem{Watrous18} J. Watrous, \emph{The Theory of Quantum Information}, Cambridge University Press, 2018, Chapter 3, Corollary 3.40.
\bibitem{KRN65} K. Kuratowski and C. Ryll-Nardzewski, \emph{A general theorem on selectors}, Bull. Acad. Polon. Sci. S\'er. Sci. Math. Astronom. Phys. \textbf{13} (1965), 397--403.

\bibitem{Charnes63} A. Charnes, W. W. Cooper and K. O. Kortanek, \emph{Duality in semi-infinite programs and some works of Haar and Carath\'eodory}, Management Sci. \textbf{9} (1963), 209--228. DOI: 10.1287/mnsc.9.2.209.
\bibitem{Shapiro09} A. Shapiro, \emph{Semi-infinite programming, duality, discretization and optimality conditions}, Optimization \textbf{58} (2009), 133--161. DOI: 10.1080/02331930902730070; author preprint, Optimization Online (2008).
\bibitem{Dufour25} F. Dufour and T. Prieto-Rumeau, \emph{Absorbing Markov decision processes and their occupation measures}, SIAM J. Control Optim. \textbf{63} (2025), 676--698. DOI: 10.1137/24M1661959. Related author preprint examined: \emph{A note on weak compactness of occupation measures for an absorbing Markov decision process}, arXiv:2402.10672v2 (2024).
\bibitem{Franc23} V. Franc, D. Prusa and V. Voracek, \emph{Optimal strategies for reject option classifiers}, J. Mach. Learn. Res. \textbf{24} (2023), no. 11, 1--49.
\bibitem{Jencova16} A. Jen\v{c}ov\'a, \emph{Comparison of quantum channels and statistical experiments}, arXiv:1512.07016v2 (2016). The separately indexed shorter version arXiv:1601.06370 was withdrawn in favor of this extended text.
\bibitem{LiElGamal18} C. T. Li and A. El Gamal, \emph{Strong functional representation lemma and applications to coding theorems}, IEEE Trans. Inform. Theory \textbf{64} (2018), 6967--6978; arXiv:1701.02827v4.
\bibitem{AldousShepp87} D. Aldous and L. Shepp, \emph{The least variable phase type distribution is Erlang}, Comm. Statist. Stochastic Models \textbf{3} (1987), 467--473. DOI: 10.1080/15326348708807067.
\bibitem{Sion58} M. Sion, \emph{On general minimax theorems}, Pacific J. Math. \textbf{8} (1958), 171--176. DOI: 10.2140/pjm.1958.8.171.
\bibitem{PiunovskiyZhang24} A. Piunovskiy and Y. Zhang, \emph{Extreme occupation measures in Markov decision processes with an absorbing state}, SIAM J. Control Optim. \textbf{62} (2024), 65--90. DOI: 10.1137/23M1572398. Author text examined: \emph{Extreme occupation measures in Markov decision processes with a cemetery}, arXiv:2307.03158v1.
\bibitem{DufourMixtures} F. Dufour and T. Prieto-Rumeau, \emph{Absorbing Markov decision processes: geometric properties and sufficiency of finite mixtures of deterministic policies}, arXiv:2512.17511; examined HTML at version-1 URL, Theorems 4.3--4.4 and Proposition 4.5.
\bibitem{Mereghetti24} C. Mereghetti, B. Palano and P. Raucci, \emph{Unary quantum finite state automata with control language}, Appl. Sci. \textbf{14} (2024), no. 4, 1490. DOI: 10.3390/app14041490.
\bibitem{RabinScott59} M. O. Rabin and D. Scott, \emph{Finite automata and their decision problems}, IBM J. Res. Develop. \textbf{3} (1959), 114--125. DOI: 10.1147/rd.32.0114.
\bibitem{Jaeger00} H. Jaeger, \emph{Observable operator models for discrete stochastic time series}, Neural Comput. \textbf{12} (2000), 1371--1398. DOI: 10.1162/089976600300015411.
\bibitem{ThonJaeger15} M. Thon and H. Jaeger, \emph{Links between multiplicity automata, observable operator models and predictive state representations---a unified learning framework}, J. Mach. Learn. Res. \textbf{16} (2015), 103--147.
\bibitem{MonrasWinter14} A. Monras and A. Winter, \emph{Quantum learning of classical stochastic processes: the completely-positive realization problem}, LIPIcs \textbf{27}, TQC 2014 (2014), 99--109. DOI: 10.4230/LIPIcs.TQC.2014.99.

\bibitem{Heller65} A. Heller, \emph{On stochastic processes derived from Markov chains}, Ann. Math. Statist. \textbf{36} (1965), 1286--1291. DOI: 10.1214/aoms/1177700000.
\bibitem{CarlylePaz71} J. W. Carlyle and A. Paz, \emph{Realizations by stochastic finite automata}, J. Comput. System Sci. \textbf{5} (1971), 26--40. DOI: 10.1016/S0022-0000(71)80005-3.
\bibitem{Fanizza24} M. Fanizza, J. Lumbreras and A. Winter, \emph{Quantum theory in finite dimension cannot explain every general process with finite memory}, Comm. Math. Phys. \textbf{405} (2024), article 50. DOI: 10.1007/s00220-023-04913-4; arXiv:2209.11225. In particular see Theorem 2 and Section 4.2 for the noise-dependent worst-phase classical-memory lower bound.
\bibitem{DmitrievDynkin46} N. Dmitriev and E. Dynkin, \emph{On characteristic roots of stochastic matrices}, Izv. Akad. Nauk SSSR Ser. Mat. \textbf{10} (1946), no. 2, 167--184. Math-Net.Ru: im3595 (Russian).
\bibitem{JohnsonPaparella17} C. R. Johnson and P. Paparella, \emph{A matricial view of the Karpelevi\v c theorem}, Linear Algebra Appl. \textbf{520} (2017), 1--15. DOI: 10.1016/j.laa.2017.01.009; arXiv:1611.06970v2.
\bibitem{KirklandSmigoc22} S. Kirkland and H. \v Smigoc, \emph{Stochastic matrices realising the boundary of the Karpelevi\v c region}, Linear Algebra Appl. \textbf{635} (2022), 116--138. DOI: 10.1016/j.laa.2021.11.016; arXiv:2110.01040v1.
\bibitem{BenvenutiFarina04} L. Benvenuti and L. Farina, \emph{A tutorial on the positive realization problem}, IEEE Trans. Automat. Control \textbf{49} (2004), 651--664. DOI: 10.1109/TAC.2004.826715.
\bibitem{BenvenutiFarina99} L. Benvenuti and L. Farina, \emph{An example of how positivity may force realizations of ``large'' dimension}, Systems Control Lett. \textbf{36} (1999), 261--266. DOI: 10.1016/S0167-6911(98)00098-X.
\bibitem{OCinneide90} C. A. O'Cinneide, \emph{Characterization of phase-type distributions}, Comm. Statist. Stochastic Models \textbf{6} (1990), 1--57. DOI: 10.1080/15326349908807134.
\bibitem{Maier91} R. S. Maier, \emph{The algebraic construction of phase-type distributions}, Comm. Statist. Stochastic Models \textbf{7} (1991), 573--602. DOI: 10.1080/15326349108807207.
\end{thebibliography}

\end{document}
